\chapter{Entwurf zeitdiskreter Systeme}\label{sec:entwurf-zeitdiskreter-systeme}

Analoge und damit zeitkontinuierliche Systeme werden in zunehmendem Maß durch zeitdiskrete Systeme ersetzt. Deshalb ist der Entwurf zeitdiskreter Systeme ein wichtiges Feld der Systemtheorie. Die zeitdiskrete Approximation zeitkontinuierlicher Systeme ist nicht eindeutig, je nach Approximationsverfahren werden unterschiedliche Ergebnisse erzielt.

Eine Möglichkeit, zeitdiskrete Systeme zu entwerfen, ergibt sich aus der zeitdiskreten Nachbildung der Systemantworten. Zum einen kann die Impulsantwort des zeitkontinuierlichen Systems zeitdiskret nachgebildet werden. Das zeitdiskrete System wird so entworfen, dass die Impulsantwort \(g[k]\) an den Stellen t = kT\textsubscript{A} dieselben Werte aufweist wie die zeitkontinuierliche Impulsantwort \(g(t)\). Alternativ kann das System so nachgebildet werden, dass die Sprungantworten von zeitdiskretem und zeitkontinuierlichem System an den Stellen t = kT\textsubscript{A} dieselben Werte aufweisen. Diese Verfahren werden als impuls- beziehungsweise sprunginvarianter Entwurf zeitdiskreter Systeme bezeichnet und im Folgenden vorgestellt.

Eine andere Möglichkeit ergibt sich aus einer Transformation der Übertragungsfunktion des zeitkontinuierlichen Systems \(G(s)\) in eine Übertragungsfunktion des zeitdiskreten Systems. Das Verfahren ist die sogenannte bilineare Transformation. Sie wird insbesondere zur Entwicklung rekursiver Filter eingesetzt. Die bilineare Transformation kann auch zur numerischen Lösung von Differentialgleichungen dynamischer Systeme eingesetzt werden.

Entwurfsverfahren für FIR-Filter unterscheiden sich von denen für IIR-Filter. Die analytischen Verfahren versuchen, ein ideales Filter bestmöglich nachzubilden. Das Verfahren ähnelt dem impulsinvarianten Entwurf für IIR-Filter. Zur Verbesserung der Filtercharakteristik können dabei Fensterfunktionen eingesetzt werden.

In der Praxis erfolgt der Entwurf von FIR-Filtern mit Hilfe von numerischen Optimierungsverfahren. Ein bekannter Vertreter der numerischen Entwurfsverfahren ist der Remez-Filterentwurf, der auf dem Remez-Austauschalgorithmus basiert.

Filter können mit unterschiedlichen Strukturen realisiert werden. Für die Realisierung rekursiver IIR-Systeme werden die Direktstruktur, die Kaskadenstruktur und die Parallelstruktur erläutert. Außerdem wird auf die Realisierung von FIR-Filtern als Kaskadenstruktur eingegangen.

\section{Nachbildung von Impuls- oder Sprungantwort}\label{sec:9.1}

\subsection{Impulsinvarianter Entwurf zeitdiskreter Systeme}\label{sec:impulsinvarianter-entwurf-zeitdiskreter-systeme}

Eine Möglichkeit der zeitdiskreten Implementierung zeitkontinuierlicher Systeme ergibt sich aus der Abtastung der Impulsantwort \(g(t)\). Technisch realisierbare zeitkontinuierliche Systeme sind kausal und stabil. Im Folgenden wird zur Vereinfachung angenommen, dass sie einfache Pole besitzen, die reell oder konjugiert komplex sein können. Da mit diesen Voraussetzungen der Zählergrad der Übertragungsfunktion bei realisierbaren Systemen kleiner oder gleich dem Nennergrad ist, lässt sich die Übertragungsfunktion als Summe einzelner Partialbrüche darstellen.
\begin{align}
G(s) &= A_{0} + \sum\limits_{n = 1}^{N}\frac{A_{n}}{s - \alpha_{n}} \label{eq:9.1}
\end{align}
Diese Darstellung führt zu der zeitkontinuierlichen Impulsantwort
\begin{align}
g(t) &= A_{0}  \delta(t) + \sum\limits_{n = 1}^{N}{A_{n}  e^{\alpha_{n}  t}  \sigma(t)} \label{eq:9.2}
\end{align}
Die zeitkontinuierliche und die zeitdiskrete Impulsantwort soll an den Stellen t = kT\textsubscript{A} dieselben Werte aufweisen. Damit ergibt sich die Impulsantwort des zeitdiskreten Systems zu
\begin{align}
g[ k] &= g\left( k  T_{A} \right) = A_{0}  \delta\left( k  T_{A} \right) + \sum\limits_{n = 1}^{N}{A_{n}  e^{\alpha_{n}  k  T_{A}}  \sigma\left( k  T_{A} \right)} \label{eq:9.3}
\end{align}
Die z-Transformierte \(G(z)\) der Impulsantwort \(g[k]\) berechnet sich über die Definitionsgleichung der\\
z-Transformation zu
\begin{align}
{G\left(z\right)} &=  {\sum\limits _{k=0}^{\infty }g\left[k\right] z^{-k}  =\sum\limits _{k=0}^{\infty }\left(A_{0}  \delta \left(k T_{A} \right)+\sum\limits _{n=1}^{N}A_{n}  e^{\alpha _{n}  k T_{A} }  \sigma \left(k T_{A} \right) \right) z^{-k}  } \\
{} &=  {A_{0} +\sum\limits _{k=0}^{\infty }\sum\limits _{n=1}^{N}A_{n}  e^{\alpha _{n}  k T_{A} }   z^{-k}  }
 \label{eq:9.4}
\end{align}
Vertauschen der Summationsreihenfolge und Anwendung der Summenformel für die geometrische Reihe führt zu
\begin{align}
{G(z) = A_{0} + \sum\limits_{n = 1}^{N}{\sum\limits_{k = 0}^{\infty}{A_{n}  e^{\alpha_{n}  k  T_{A}}  z^{- k}}} = A_{0} + \sum\limits_{n = 1}^{N}{A_{n}  \sum\limits_{k = 0}^{\infty}\left( e^{\alpha_{n}  T_{A}}  z^{- 1} \right)^{k}}
}{= A_{0} + \sum\limits_{n = 1}^{N}{A_{n}  \frac{1}{1 - e^{\alpha_{n}  T_{A}}  z^{- 1}}} = A_{0} + \sum\limits_{n = 1}^{N}{A_{n}  \frac{z}{z - e^{\alpha_{n}  T_{A}}}}} \label{eq:9.5}
\end{align}
Diese Übertragungsfunktion hat einen Konvergenzbereich von
\begin{align}
\left| e^{\alpha_{n}  T_{A}} \right| < |z| \label{eq:9.6}
\end{align}
Da die Polstellen \(α_{n}\) zu einem stabilen zeitkontinuierlichen System gehören, liegen sie in der negativen Halbebene. Damit liegen die Pole der Übertragungsfunktion \(G(z)\) innerhalb des Einheitskreises. Beim impulsinvarianten Entwurf bleiben stabile Systeme demnach stabil.
\begin{application}{Transformation}
  Im Online-Portal \emph{Systemtheorie Online} verdeutlicht die Applikation \emph{Transformation} den Zusammenhang zwischen Punkten in der Laplace-Ebene und Punkten in der z-Ebene für unterschiedliche Transformationsverfahren.
  \tcblower
  \url{h-ka.de/eit/digitale-signalverarbeitung/übungen/kapitel2}
\end{application}
\subparagraph{Beispiel: Impulsinvarianter Entwurf eines Tiefpasses erster Ordnung}\label{beispiel-impulsinvarianter-entwurf-eines-tiefpasses-erster-ordnung}

Um das Verfahren zu verdeutlichen, wird es an einem Tiefpass erster Ordnung angewendet. Das Übertragungsglied besitzt die Übertragungsfunktion
\begin{align}
G(s) &= \frac{\omega_{G}}{s + \omega_{G}} \label{eq:9.7}
\end{align}
Es weist eine Verstärkung von G(\(s = 0\)) = 1 auf und die zeitkontinuierliche Impulsantwort lauet
\begin{align}
g(t) &= \omega_{G}  e^{\omega_{G}  t}  \sigma(t) \label{eq:9.8}
\end{align}
Mit dem Pol α = - \(ω_{G}\) ergibt sich die Übertragungsfunktion des zeitdiskreten Systems von
\begin{align}
G(z) &= \frac{Y(z)}{U(z)} = \omega_{G}  \frac{z}{z - e^{- \omega_{G}  T_{A}}} = \omega_{G}  \frac{1}{1 - e^{- \omega_{G}  T_{A}}  z^{- 1}} \label{eq:9.9}
\end{align}
Das System kann mit der Differenzengleichung
\begin{align}
y[ k] &= \omega_{G}  u[ k] + e^{- \omega_{G}  T_{A}}  y[ k - 1] \label{eq:9.10}
\end{align}
realisiert werden und hat eine Verstärkung von G(\(z = 1\)) = 1.103. Bild \ref{fig:9.1} vergleicht das zeitkontinuierliche und das zeitdiskrete System im Zeitbereich für \(T = 5\) und \(T_{A} = 1\).

\begin{figure}
  \centering
  \includegraphics[alt={Vergleich des zeitkontinuierlichen und des zeitdiskreten Systems\\
im Zeitbereich bei impulsinvariantem Entwurf für \(T = 5\) und \(T_{A} = 1\)}]{new media/10_1_vergleich_des_zeitkontinuierlic.png}
  \caption{Vergleich des zeitkontinuierlichen und des zeitdiskreten Systems\\
im Zeitbereich bei impulsinvariantem Entwurf für \(T = 5\) und \(T_{A} = 1\)}
  \label{fig:9.1}
\end{figure}

Die beiden Impulsantworten \(g(t)\) und \(g[k]\) stimmen erwartungsgemäß an den Stellen t = kT\textsubscript{A} überein. Die Sprungantworten weichen in ihrem Zeitverlauf und in ihrem Endwert voneinander ab. Bild \ref{fig:9.2} vergleicht die beiden Systeme im Frequenzbereich.

\begin{figure}
  \centering
  \includegraphics[alt={Vergleich des zeitkontinuierlichen und des zeitdiskreten Systems\\
im Frequenzbereich bei impulsinvariantem Entwurf für \(T = 5\) und \(T_{A} = 1\)}]{new media/10_2_vergleich_des_zeitkontinuierlic.png}
  \caption{Vergleich des zeitkontinuierlichen und des zeitdiskreten Systems\\
im Frequenzbereich bei impulsinvariantem Entwurf für \(T = 5\) und \(T_{A} = 1\)}
  \label{fig:9.2}
\end{figure}

Die beiden Amplitudengänge stimmen, abgesehen von dem unterschiedlichen Verstärkungsfaktor, im Bereich kleiner Frequenzen sehr gut überein. Die Abweichungen im Frequenzbereich π/2 \textless{} Ω \textless{} π sind auf die periodische Wiederholung des Spektrums zurückzuführen. Da der Amplitudengang an der Stelle \(Ω = π\) noch nicht vollständig verschwindet, wird die periodische Fortsetzung in das Basisband gefaltet. Dieser Effekt sinkt mit steigender Zeitkonstante T des Systems, da damit die eigentlich erforderliche Bandbegrenzung des Systems besser realisiert wird. Alternativ kann die Abtastzeit reduziert werden.

Die beiden Phasengänge stimmen nur im Bereich sehr kleiner Frequenzen überein. Eigentliche Ursache für die hohen Abweichungen im Phasengang sind die durch den Abtastvorgang entstehenden zusätzlichen Nullstellen in der Übertragungsfunktion. Bild \ref{fig:9.3} vergleicht die Pol-Nullstellen-Diagramme des zeitkontinuierlichen und des zeitdiskreten Systems.

\begin{figure}
  \centering
  \includegraphics[alt={Vergleich der Pol-Nullstellen-Diagramme}]{new media/10_3_vergleich_der_polnullstellendia.png}
  \caption{Vergleich der Pol-Nullstellen-Diagramme}
  \label{fig:9.3}
\end{figure}

Beide Systeme weisen einen einfachen reellen Pol auf. Die zusätzliche Nullstelle im zeitdiskreten System führt bei normierten Frequenzen Ω = ± π zu einer Phase ϕ = 0.

▪

Das Vorgehen zum impulsinvarianten Entwurf zeitdiskreter Systeme ist in Tabelle \ref{tab:9.1} zusammengefasst.

\begin{table}
  \centering
  \caption{Vorgehen zum impulsinvarianten Entwurf zeitdiskreter Systeme bei einfachen Polen \(α_{n}\)}
  \label{tab:9.1}
  \begin{tabular}{ll}
    \toprule
    Schritt & Beschreibung \\
    \midrule
    1 & Entwurf eines zeitkontinuierlichen Systems \(G(s)\) \\
    2 & Darstellung des Systems als Summe von Partialbrüchen \\
    & \(G(s) = A_{0} + \sum\limits_{n = 1}^{N}\frac{A_{n}}{s - \alpha_{n}}\) \\
    3 & Zeitdiskretes System besitzt die Übertragungsfunktion \(G(z)\) \\
    & \(G(z) = A_{0} + \sum\limits_{n = 1}^{N}{A_{n}  \frac{1}{1 - e^{\alpha_{n}  T_{A}}  z^{- 1}}}\) \\
    4 & Aufstellen der Differenzengleichung durch Rücktransformation der Gleichung \(G(z)\) in den Zeitbereich \\
    & \(G(z) = \frac{Y(z)}{U(z)}\) \\
    5 & Auflösen der Differenzengleichung nach \(y[k]\) \\
    \bottomrule
  \end{tabular}
\end{table}

Die Transformation der Laplace-Transformierten in den Zeitbereich, die Abtastung der Zeitfunktion und die anschließende Transformation in den z-Bereich wird auch als zusammengesetzte z-Transformation bezeichnet. Sie wir symbolisch dargestellt als
\begin{align}
G(z) &= Z\left\{ G(s) \right\} \label{eq:9.11}
\end{align}
Der impulsinvariante Entwurf wird in der Regelungstechnik für die zeitdiskrete Nachbildung analoger Systeme verwendet. Liegt zum Beispiel ein analoger Reglerentwurf vor, kann der zeitkontinuierliche Regler in einen zeitdiskreten Regler überführt werden. Zeitdiskrete Filter lassen sich entwickeln, in dem ein zeitkontinuierliches Filter mithilfe des impulsinvarianten Entwurfs in ein zeitdiskretes Filter überführt wird.

\subsection{Sprunginvarianter Entwurf zeitdiskreter Systeme}\label{sec:sprunginvarianter-entwurf-zeitdiskreter-systeme}

Eine weitere Möglichkeit der zeitdiskreten Implementierung zeitkontinuierlicher Systeme ergibt sich aus dem Ziel, die Sprungantwort nachzubilden. Ausgangspunkt für die Herleitung ist wieder die Darstellung der Übertragungsfunktion in Partialbrüchen. Wieder wird zur Vereinfachung angenommen, dass die Übertragungsfunktion einfache reelle Pole und konjugiert komplexe Polpaare besitzt.
\begin{align}
G(s) &= A_{0} + \sum\limits_{n = 1}^{N}\frac{A_{n}}{s - \alpha_{n}} \label{eq:9.12}
\end{align}
Die Sprungantwort hat damit die Laplace-Transformierte
\begin{align}
H(s) &= \frac{1}{s}  G(s) = \frac{A_{0}}{s} + \sum\limits_{n = 1}^{N}{\frac{1}{s}  \frac{A_{n}}{s - \alpha_{n}}} = \frac{A_{0}}{s} + \sum\limits_{n = 1}^{N}{\frac{A_{n}}{- \alpha_{n}}  \frac{1}{s}  \frac{1}{1 - \frac{1}{\alpha_{n}}  s}} \label{eq:9.13}
\end{align}
Damit lautet die Sprungantwort im Zeitbereich
\begin{align}
h(t) &= A_{0}  \sigma(t) + \sum\limits_{n = 1}^{N}{\frac{A_{n}}{- \alpha_{n}}  \left( 1 - e^{\alpha_{n}  t} \right)  \sigma(t)} = \left( A_{0} - \sum\limits_{n = 1}^{N}\frac{A_{n}}{\alpha_{n}} \right)  \sigma(t) + \sum\limits_{n = 1}^{N}{\frac{A_{n}}{\alpha_{n}}  e^{\alpha_{n}  t}  \sigma(t)} \label{eq:9.14}
\end{align}
An den Zeitpunkten kT\textsubscript{A} liegen die Werte
\begin{align}
h\left( k  T_{A} \right) &= \left( A_{0} - \sum\limits_{n = 1}^{N}\frac{A_{n}}{\alpha_{n}} \right)  \sigma\left( k  T_{A} \right) + \sum\limits_{n = 1}^{N}{\frac{A_{n}}{\alpha_{n}}  e^{\alpha_{n}  k  T_{A}}  \sigma\left( k  T_{A} \right)} \label{eq:9.15}
\end{align}
vor. Damit lautet die Sprungantwort des zeitdiskret nachzubildenden Systems
\begin{align}
h[ k] &= \left( A_{0} - \sum\limits_{n = 1}^{N}\frac{A_{n}}{\alpha_{n}} \right)  \sigma[ k] + \sum\limits_{n = 1}^{N}{\frac{A_{n}}{\alpha_{n}}  e^{\alpha_{n}  k  T_{A}}  \sigma[ k]} \label{eq:9.16}
\end{align}
Die Sprungantwort hat die z-Transformierte
\begin{align}
H(z) &= \left( A_{0} - \sum\limits_{n = 1}^{N}\frac{A_{n}}{\alpha_{n}} \right)  \frac{z}{z - 1} + \sum\limits_{n = 1}^{N}{\frac{A_{n}}{\alpha_{n}}  \frac{z}{z - e^{\alpha_{n}  T_{A}}}} \label{eq:9.17}
\end{align}
Aus ihr ergibt sich die Übertragungsfunktion des zeitdiskreten Systems zu
\begin{align}
{G\left(z\right)} &=  {\frac{z-1}{z}  H\left(z\right)=\left(A_{0} -\sum\limits _{n=1}^{N}\frac{A_{n} }{\alpha _{n} }  \right) \frac{z}{z-1}  \frac{z-1}{z} +\sum\limits _{n=1}^{N}\frac{A_{n} }{\alpha _{n} }  \frac{z}{z-e^{\alpha _{n}  T_{A} } }  \frac{z-1}{z}  } \\
{} &=  {\left(A_{0} -\sum\limits _{n=1}^{N}\frac{A_{n} }{\alpha _{n} }  \right)+\sum\limits _{n=1}^{N}\frac{A_{n} }{\alpha _{n} }  \frac{z-1}{z-e^{\alpha _{n}  T_{A} } }  }
 \label{eq:9.18}
\end{align}
Diese Übertragungsfunktion hat einen Konvergenzbereich von
\begin{align}
\left| e^{\alpha_{n}  T_{A}} \right| < |z| \label{eq:9.19}
\end{align}
Da die Polstellen \(α_{n}\) zu einem stabilen zeitkontinuierlichen System gehören, liegen sie in der negativen Halbebene. Damit liegen die Pole der Übertragungsfunktion \(G(z)\) innerhalb des Einheitskreises. Beim sprunginvarianten Entwurf bleiben stabile Systeme demnach stabil.

\begin{example}{Sprunginvarianter Entwurf eines Tiefpasses erster Ordnung}
\label{beispiel-sprunginvarianter-entwurf-eines-tiefpasses-erster-ordnung}

Der sprunginvariante Entwurf wird ebenfalls an einem Tiefpass erster Ordnung durchgeführt.
\begin{align}
G(s) &= \frac{1}{1 + T  s} = \frac{\omega_{G}}{s + \omega_{G}} \label{eq:9.20}
\end{align}
Mit dem Pol α = - \(ω_{G}\) ergibt sich beim sprunginvarianten Entwurf die Übertragungsfunktion des zeitdiskreten Systems zu
\begin{align}
{G\left(z\right)} &=  {-\frac{\omega _{G} }{-\omega _{G} } +\frac{\omega _{G} }{-\omega _{G} }  \frac{z-1}{z-e^{-\omega _{G}  T_{A} } } =1-\frac{z-1}{z-e^{-\omega _{G}  T_{A} } } } \\
{} &=  {\frac{z-e^{-\omega _{G}  T_{A} } -z+1}{z-e^{-\omega _{G}  T_{A} } } =\frac{1-e^{-\omega _{G}  T_{A} } }{z-e^{-\omega _{G}  T_{A} } } }
 \label{eq:9.21}
\end{align}
Bild \ref{fig:9.1} vergleicht das zeitkontinuierliche und das zeitdiskrete System im Zeitbereich für \(T = 5\) und\\
\(T_{A} = 1\).

\begin{fixedfigure}
  \includegraphics[alt={Vergleich der Impuls- und Sprungantworten von zeitdiskretem und zeitkontinuierlichem System}]{new media/10_4_vergleich_der_impuls_und_sprung.png}
  \captionof{figure}{Vergleich der Impuls- und Sprungantworten von zeitdiskretem und zeitkontinuierlichem System}
\end{fixedfigure}

Im Gegensatz zum impulsinvarianten Entwurf stimmen beim sprunginvarianten Entwurf die Sprungantworten überein.

\end{example}

Das Vorgehen zum sprunginvarianten Entwurf zeitdiskreter Systeme ist in Tabelle \ref{tab:9.2} zusammengefasst.

\begin{table}
  \centering
  \caption{Vorgehen zum sprunginvarianten Entwurf zeitdiskreter Systeme bei einfachen Polen αn ≠ 0}
  \label{tab:9.2}
  \begin{tabular}{ll}
    \toprule
    Schritt & Beschreibung \\
    \midrule
    1 & Entwurf eines zeitkontinuierlichen Systems \(G(s)\) \\
    2 & Darstellung des Systems als Summe von Partialbrüchen \\
    & \(G(s) = A_{0} + \sum\limits_{n = 1}^{N}\frac{A_{n}}{s - \alpha_{n}}\) \\
    3 & Zeitdiskretes System besitzt die Übertragungsfunktion \(G(z)\) \\
    & \(G(z) = \left( A_{0} - \sum\limits_{n = 1}^{N}\frac{A_{n}}{\alpha_{n}} \right) + \sum\limits_{n = 1}^{N}{\frac{A_{n}}{\alpha_{n}}  \frac{z - 1}{z - e^{- \alpha_{n}  T_{A}}}}\) \\
    4 & Aufstellen der Differenzengleichung durch Rücktransformation der Gleichung \(G(z)\) in den Zeitbereich \\
    & \(G(z) = \frac{Y(z)}{U(z)}\) \\
    5 & Auflösen der Differenzengleichung nach \(y[k]\) \\
    \bottomrule
  \end{tabular}
\end{table}

Das Verfahren wird auch als \emph{Zero-Order-Hold-Equivalence} bezeichnet. Die Bezeichnung ergibt sich aus dem Ansatz, ein zeitkontinuierliches System \(G(s)\) mit zeitdiskreter Ansteuerung \(u[k]\) zeitdiskret zu simulieren. Bild \ref{fig:9.5} stellt den Signalfluss dar.

\begin{figure}
  \centering
  \includesvg{new media/10_5_streckendiskretisierung_mit_ein.svg}
  \caption{Strecken-Diskretisierung mit einem Zero-Order-Hold-Glied und einer idealen Abtastung}
  \label{fig:9.5}
\end{figure}

Um aus dem zeitdiskreten Wert ein zeitkontinuierliches Signal zu bekommen, wird das Signal konstant gehalten, bis der nächste Abtastwert vorliegt. Dieser Vorgang wird mit einem Zero-Order-Hold-Glied modelliert, das Ausgangssignal des Zero-Order-Hold-Glieds wird als \(u_{H}(t)\) bezeichnet.
\begin{align}
u_{H}(t) &= u_{A}(t)*\left( \sigma(t) - \sigma\left( t - T_{A} \right) \right) \label{eq:9.22}
\end{align}
Im Laplace-Bereich hat das Halteglied die Übertragungsfunktion
\begin{align}
G_{H}(s) &= \frac{1 - e^{- s  T_{A}}}{s} \label{eq:9.23}
\end{align}
Das zeitkontinuierliche Signal \(u_{H}(t)\) ist das Eingangssignal des zeitkontinuierlichen Systems mit der Übertragungsfunktion \(G(s)\). Das entsprechende Ausgangssignal \(y(t)\) wird ideal abgetastet. Es entsteht die Abtastfolge \(y[k]\). Die dabei entstehende Übertragungsfunktion \(G(z)\) der abgetasteten Strecke berechnet sich zu
\begin{align}
G(z) &= Z\left\{ \frac{1 - e^{- s  T_{A}}}{s}  G(s) \right\} = \left( 1 - z^{- 1} \right)  Z\left\{ \frac{1}{s}  G(s) \right\} = \frac{z - 1}{z}  Z\left\{ G(s)  \frac{1}{s} \right\} \label{eq:9.24}
\end{align}
Umformen führt zu
\begin{align}
\frac{z}{z - 1}  G(z) &= Z\left\{ G(s)  \frac{1}{s} \right\} \label{eq:9.25}
\end{align}
beziehungsweise
\begin{align}
H(z) &= Z\left\{ H(s) \right\} \label{eq:9.26}
\end{align}
Die Sprungantworten stimmen damit an den Abtastzeitpunkten überein, es handelt sich um einen sprunginvarianten Entwurf.

\subsection{Korrespondenzen für den impuls- und sprunginvarianten Entwurf}\label{sec:korrespondenzen-fuxfcr-den-impuls--und-sprunginvarianten-entwurf}

In Korrespondenztafeln kann der Zusammenhang zwischen Übertragungsfunktionen im Laplace- und z-Bereich direkt nachgeschlagen werden.

\begin{table}
  \centering
  \caption{Korrespondenzen für den impuls- und sprunginvarianten Entwurf}
  \begin{tabular}{llll}
    \toprule
    Nr. & 
Laplace- 
Transformierte 
\(G(s)\)
 & \multicolumn{2}{l}{%
z-Transformierte \(G(z)\)} \\
    \midrule
    & & 
Impulsinvarianter 
Entwurf
 & 
Sprunginvarianter 
Entwurf \\
    1 & \(\frac{1}{s}\) & \(\frac{z}{z - 1}\) & \(\frac{T_{A}}{z - 1}\) \\
    2 & \(\frac{1}{s^{2}}\) & \(\frac{T_{A}  z}{(z - 1)^{2}}\) & \(\frac{T_{A}^{2}}{2}  \frac{z + 1}{(z - 1)^{2}}\) \\
    3 & \(\frac{1}{s - \alpha}\) & \(\frac{z}{z - e^{\alpha  T_{A}}}\) & \(\frac{1 - e^{\alpha  T_{A}}}{- \alpha  \left( z - e^{\alpha  T_{A}} \right)}\) \\
    4 & \(\frac{1}{(s - \alpha)^{2}}\) & \(\frac{T_{A}  e^{\alpha  T_{A}}  z}{\left( z - e^{\alpha  T_{A}} \right)^{2}}\) & \(\frac{\left( 1 - e^{\alpha  T_{A}} + \alpha  T_{A}  e^{\alpha  T_{A}} \right)  z + \left( e^{\alpha  T_{A}} - \alpha  T_{A} - 1 \right)  e^{\alpha  T_{A}}}{\alpha^{2}  \left( z - e^{\alpha  T_{A}} \right)^{2}}\) \\
    \bottomrule
  \end{tabular}
\end{table}

\subsection{Impulsinvarianter Filterentwurf eines IIR-Filters}\label{sec:9.1.4}

In Teil A dieser Buchreihe wird ein Butterworth-Filter 3. Ordnung mit einer Grenzfrequenz\\
\(f_{G} = \SI{10}{\kilo\hertz}\) und einer maximalen Durchlassdämpfung von - \SI{3}{\decibel} entworfen. Das System besitzt die Übertragungsfunktion
\begin{align}
{G\left(s\right)} &=  {\frac{\omega _{G}^{3} }{\left(s+\omega _{G} \right) \left(s+\left(\frac{1}{2} -j \frac{\sqrt{3} }{2} \right) \omega _{G} \right) \left(s+\left(\frac{1}{2} +j \frac{\sqrt{3} }{2} \right) \omega _{G} \right)} } \\
{} &=  {\frac{\omega _{G} }{s+\omega _{G} } +\frac{\frac{2}{-3+j \sqrt{3} }  \omega _{G} }{s+\left(\frac{1}{2} -j \frac{\sqrt{3} }{2} \right) \omega _{G} } +\frac{\frac{2}{-3-j \sqrt{3} }  \omega _{G} }{s+\left(\frac{1}{2} +j \frac{\sqrt{3} }{2} \right) \omega _{G} } }
 \label{eq:9.27}
\end{align}
Das Filter soll zeitdiskret auf einem System umgesetzt werden, das mit einer Abtastrate von\\
\(f_{A} = \SI{100}{\kilo\hertz}\) arbeitet. Nach Gleichung \eqref{eq:9.5} ergibt sich die z-Transformierte \(G(z)\) zu
\begin{align}
G(z) &= \frac{\omega_{G}}{1 - e^{- T_{A}  \omega_{G}}  z^{- 1}} + \frac{\omega_{G}  \frac{2}{- 3 + j  \sqrt{3}}}{1 - e^{- \left( \frac{1}{2} - j  \frac{\sqrt{3}}{2} \right)  T_{A}  \omega_{G}}  z^{- 1}} + \frac{\omega_{G}  \frac{2}{- 3 - j  \sqrt{3}}}{1 - e^{- \left( \frac{1}{2} + j  \frac{\sqrt{3}}{2} \right)  T_{A}  \omega_{G}}  z^{- 1}} \label{eq:9.28}
\end{align}
Mit der indirekt angegebenen Abtastzeit von \(T_{A}\) = \(10^{-5}\) s und der Grenzfrequenz \(ω_{G}\) = 2π\SI{10}{\kilo\hertz} ergibt sich das Produkt
\begin{align}
T_{A}  \omega_{G} &= 10^{- 5}\text{ s}  2  \pi  10^{4}\ \frac{1}{s} = \frac{2  \pi}{10} = \frac{\pi}{5} \label{eq:9.29}
\end{align}
und die Übertragungsfunktion \(G(z)\) berechnet sich zu
\begin{align}
{G\left(z\right)} &=  {\frac{\omega _{G} }{1-e^{-\frac{\pi }{5} }  z^{-1} } +\frac{\omega _{G}  \frac{2}{-3+j \sqrt{3} } }{1-e^{-\left(\frac{1}{2} -j \frac{\sqrt{3} }{2} \right) \frac{\pi }{5} }  z^{-1} } +\frac{\omega _{G}  \frac{2}{-3-j \sqrt{3} } }{1-e^{-\left(\frac{1}{2} +j \frac{\sqrt{3} }{2} \right) \frac{\pi }{5} }  z^{-1} } } \\
{} &=  {{79734} \frac{z^{2} +0.6587 z}{z^{3} -1.7833 z^{2} +1.2002 z-0.2846} }
 \label{eq:9.30}
\end{align}
Bild \ref{fig:9.6} stellt das Pol-Nullstellendiagramm für das impulsinvariant entworfene Butterworth-Filter dar.

\begin{figure}
  \centering
  \includegraphics[alt={Pol-Nullstellendiagramm für das impulsinvariant entworfene Butterworth-Filter dritter Ordnung}]{new media/10_6_polnullstellendiagramm_fr_das_i.png}
  \caption{Pol-Nullstellendiagramm für das impulsinvariant entworfene Butterworth-Filter dritter Ordnung}
  \label{fig:9.6}
\end{figure}

Die Pole liegen erwartungsgemäß im Inneren des Einheitskreises, das System ist demnach stabil. Die Verstärkung des zeitdiskreten Systems ist im Gegensatz zu dem Verstärkungsfaktor des zeitkontinuierlichen Systems ungleich eins. Er errechnet sich zu
\begin{align}
G(1) &= \text{79734}  \frac{1 + 0.6587}{1 - 1.7833 + 1.2002 - 0.2846} = \text{999650} \label{eq:9.31}
\end{align}
Bei einem direkten Vergleich der beiden Systeme muss mit diesem Faktor entsprechend normiert werden. Der Frequenzgang \(G(Ω)\) berechnet sich zu
\begin{align}
G(\Omega) &= 0.07976  \frac{e^{j  2  \Omega} + 0.6587  e^{j  \Omega}}{e^{j  3  \Omega} - 1.7833  e^{j  2  \Omega} + 1.2002  e^{j  \Omega} - 0.2846} \label{eq:9.32}
\end{align}
Bild \ref{fig:9.7} zeigt den normierten Frequenzgang des zeitdiskreten Systems im linearen Maßstab im Vergleich zu dem Frequenzgang des zeitkontinuierlichen Systems.

\begin{figure}
  \centering
  \includegraphics[alt={Frequenzgang des impulsinvariant entworfenen Butterworth-Filters dritter Ordnung}]{new media/10_7_frequenzgang_des_impulsinvarian.png}
  \caption{Frequenzgang des impulsinvariant entworfenen Butterworth-Filters dritter Ordnung}
  \label{fig:9.7}
\end{figure}

Der Amplitudengang ist für das zeitdiskrete und das zeitkontinuierliche Filter identisch. Die\\
\SI{3}{\decibel}-Grenzfrequenz des zeitdiskreten Filters beträgt \SI{10}{\kilo\hertz} und stimmt mit der 3-dB-Grenzgfrequenz des zeitkontinuierlichen Filters überein. Da der Amplitudengang für \(Ω = π\) nahe null ist, wirkt sich die periodische Wiederholung nicht auf das Basisband aus. Der Phasengang stimmt bis zu der Frequenz\\
\(Ω = π\)/2 gut überein. Bei größeren Frequenzen führt die periodische Wiederholung des Phasengangs zu einem Phasenanstieg.

Mit der Übertragungsfunktion \(G(z)\) kann das Filter als Differenzengleichung realisiert werden. Ausgehend von der Übertragungsfunktion ergibt sich nach der Division von Zähler und Nenner durch die höchste Potenz von z die Gleichung
\begin{align}
G(z) &= \frac{Y(z)}{U(z)} = 0.07976  \frac{z^{- 1} + \text{0.6587}  z^{- 2}}{1 - 1.7833  z^{- 1} + 1.2002  z^{- 2} - 0.2846  z^{- 3}} \label{eq:9.33}
\end{align}
und damit die Differenzengleichung
\begin{align}
{y\left[k\right]} &=  {1.7833 y\left[k-1\right]-1.2002 y\left[k-2\right]+0.2846 y\left[k-3\right]} \\
{} & {} & {+0.07976 u\left[k-1\right]+{0.05254} u\left[k-2\right]}
 \label{eq:9.34}
\end{align}
Bild \ref{fig:9.8} stellt die Impuls- und Sprungantwort des zeitdiskreten Systems im Vergleich zum zeitkontinuierlichen System dar. Wieder werden die zeitdiskreten Signale entsprechend normiert.

\begin{figure}
  \centering
  \includegraphics[alt={Impuls- und Sprungantwort für das impulsinvariant entworfene Butterworth-Filter dritter Ordnung}]{new media/10_8_impuls_und_sprungantwort_fr_das.png}
  \caption{Impuls- und Sprungantwort für das impulsinvariant entworfene Butterworth-Filter dritter Ordnung}
  \label{fig:9.8}
\end{figure}

Charakteristisch für den impulsvarianten Entwurf ist das Übereinstimmen der zeitkontinuierlichen und der zeitdiskreten Impulsantwort. Da die Sprungantwort im zeitkontinuierlichen Bereich über ein Integral und im zeitdiskreten Bereich über eine Summe gebildet wird, sind sie trotz identischer Impulsantwort unterschiedlich.

\section{Approximation des Systemverhaltens}\label{sec:9.2}

Bei der Einführung der z-Transformation wird die zeitliche Verschiebung um \(T_{A}\) mit der Variablen z beschrieben.
\begin{align}
z &= e^{T_{A}  s} \label{eq:9.35}
\end{align}
Dadurch wird einen Punkt der komplexen s-Ebene ein Punkt der komplexen z-Ebene zugeordnet. Um von der gebrochen rationalen Übertragungsfunktion \(G(s)\) zu einer Übertragungsfunktion \(G(z)\) zu kommen, müsste die Variable s substituiert werden durch
\begin{align}
s &= \frac{1}{T_{A}}  \ln(z) \label{eq:9.36}
\end{align}
Damit wäre die Übertragungsfunktion \(G(z)\) aber keine rationale Funktion mehr. Die Realisierung des Systems mit einer Differenzengleichung wäre nicht möglich. Es muss demnach eine Abbildungsvorschrift eingesetzt werden, die eine gebrochen rationale Funktion in s in eine gebrochen rationale Funktion in z überführt. Ausgangspunkt für die Herleitung der Abbildungsvorschrift ist ein Integrator, der für kausale Signale zeitkontinuierlich mit der Gleichung
\begin{align}
y(t) &= \int_{0}^{t}{u(\tau)}\ d\tau \label{eq:9.37}
\end{align}
beschrieben wird. Im Laplace-Bereich besitzt er die Übertragungsfunktion
\begin{align}
G(s) &= \frac{Y(s)}{U(s)} = \frac{1}{s} \label{eq:9.38}
\end{align}
Zur Herleitung der Transformationsvorschrift vom Laplace-Bereich in den z-Bereich werden verschiedene zeitdiskrete Verfahren untersucht, mit denen der Integrator approximiert wird. Im Einzelnen sind das die Verfahren

\begin{itemize}
\item
  Backward-Euler-Verfahren
\item
  Forward-Euler-Verfahren
\item
  Trapezregel oder bilineare Transformation
\end{itemize}

Bild \ref{fig:9.9} vergleicht die Grundidee der unterschiedlichen Verfahren.

\begin{figure}
  \centering
  \includegraphics[alt={Unterschiedliche Verfahren zur zeitdiskreten Approximation eines Integrators}]{new media/10_9_unterschiedliche_verfahren_zur.png}
  \caption{Unterschiedliche Verfahren zur zeitdiskreten Approximation eines Integrators}
  \label{fig:9.9}
\end{figure}

Das Integral einer Funktion entspricht der Fläche unter der Kurve. Sind von der zu integrierenden Kurve nur die Stützstellen u(kT\textsubscript{A}) = \(u[k]\) bekannt, kann die Fläche mit diesen Punkten nur approximiert werden.
\begin{application}{Transformation}
  Im Online-Portal \emph{Systemtheorie Online} verdeutlicht die Applikation \emph{Transformation} den Zusammenhang zwischen Punkten in der Laplace-Ebene und Punkten in der z-Ebene für unterschiedliche Transformationsverfahren.
  \tcblower
  \url{h-ka.de/eit/digitale-signalverarbeitung/übungen/kapitel2}
\end{application}
\subsection{Backward-Euler-Verfahren}\label{sec:backward-euler-verfahren}

Das Backward-Euler-Verfahren verwendet zur Approximation des Integrals die Gleichung
\begin{align}
y[ k] &= y[ k - 1] + u[ k]  T_{A} \label{eq:9.39}
\end{align}
Der aktuelle Wert des Ausgangsignals ergibt sich aus dem letzten Integrationswert \(y[k - 1]\) und der Approximation der Fläche in dem neuen Segment über ein Rechteck der Breite \(T_{A}\) und der Höhe\\
\(u[k]\). Die Transformation der Gleichung in den z-Bereich führt zu der Übertragungsfunktion
\begin{align}
G(z) &= \frac{Y(z)}{U(z)} = \frac{T_{A}}{1 - z^{- 1}} = \frac{z  T_{A}}{z - 1} \label{eq:9.40}
\end{align}
Ein Vergleich mit der Übertragungsfunktion des Integrators im Laplace-Bereich nach Gleichung \eqref{eq:9.38} führt zu der Transformation der Variable s des Laplace-Bereiches durch die Transformationsgleichung
\begin{align}
s &= \frac{z - 1}{z  T_{A}} \label{eq:9.41}
\end{align}
Umgekehrt kann die Variable z des zeitdiskreten Systems ersetzt werden durch
\begin{align}
z &= \frac{1}{1 - s  T_{A}} \label{eq:9.42}
\end{align}
\begin{example}{Entwurf eines Tiefpasses erster Ordnung nach dem Backward-Euler-Verfahren}
\label{beispiel-entwurf-eines-tiefpasses-erster-ordnung-nach-dem-backward-euler-verfahren}

Für ein PT1-Glied ergibt sich mit dieser Transformation die Übertragungsfunktion
\begin{align}
G(z) &= \frac{Y(z)}{U(z)} = \left. \ G(s) \right|_{s = \frac{z - 1}{z  T_{A}}} = \frac{1}{1 + T  \frac{z - 1}{z  T_{A}}} = \frac{z  T_{A}}{z  T_{A} + T  z - T} = \frac{T_{A}}{T_{A} + T - T  z^{- 1}} \label{eq:9.43}
\end{align}
Das System könnte mit der Differenzengleichung
\begin{align}
y[ k] &= \frac{T_{A}  u[ k] + T  y[ k - 1]}{T_{A} + T} \label{eq:9.44}
\end{align}
implementiert werden.

\end{example}

Der Zusammenhang zwischen Laplace- und z-Transformation wird in Kapitel \ref{sec:5.1.4} diskutiert. Der dort herausgearbeitete Zusammenhang
\begin{align}
z &= e^{s  T_{A}} \label{eq:9.45}
\end{align}
wird im Folgenden mit Zusammenhang
\begin{align}
z &= \frac{1}{1 - s  T_{A}} \label{eq:9.46}
\end{align}
verglichen, der sich aus dem Backward-Euler-Verfahren ergibt. Durch Einsetzen der Bedingung\\
s = jω kann die Abbildung der imaginären Achse der s-Ebene auf die z-Ebene bestimmt werden.
\begin{align}
z &= \frac{1}{1 - j  \omega  T_{A}} = \frac{1}{1 - j  \Omega} = \frac{1}{\sqrt{1 + \Omega^{2}}}  e^{- j  \arctan(\Omega)} \label{eq:9.47}
\end{align}
Bild \ref{fig:9.10} stellt die Abbildung der imaginären Achse der s-Ebene auf die z-Ebene bei Verwendung des Backward-Euler-Verfahrens dar.

\begin{figure}
  \centering
  \includegraphics[alt={Abbildung der imaginären Achse der s-Ebene auf die z-Ebene\\
mit dem Backward-Euler-Verfahren}]{new media/10_10_abbildung_der_imaginren_achse.png}
  \caption{Abbildung der imaginären Achse der s-Ebene auf die z-Ebene\\
mit dem Backward-Euler-Verfahren}
  \label{fig:9.10}
\end{figure}

Nach der Definitionsgleichung der z-Transformation wird die imaginäre Achse der Laplace-Ebene auf den Einheitskreis der z-Ebene abgebildet. Mit dem Backward-Euler-Verfahren wird die imaginäre Achse der Laplace-Ebene im Gegensatz dazu auf einen Kreisausschnitt in der positiven Halbebene abgebildet. Die negative s-Halbebene liegt in beiden Fällen innerhalb des jeweiligen Kreises. Die beiden Transformationen führen damit zu unterschiedlichen Abbildungen der Variable s = jω in die\\
z-Ebene. Bei dem Backward-Euler-Verfahren wird ein stabiles zeitkontinuierliches System \(G(s)\) in ein stabiles zeitdiskretes System \(G(z)\) abgebildet. Dies gilt allerdings nicht für die Rücktransformation. Ein stabiles zeitdiskretes System wird nicht unbedingt in ein zeitkontinuierliches stabiles abgebildet.

\subsection{Forward-Euler-Verfahren}\label{sec:forward-euler-verfahren}

Das Forward-Euler-Verfahren verwendet zur Approximation des Integrals die Gleichung
\begin{align}
y[ k] &= y[ k - 1] + u[ k - 1]  T_{A} \label{eq:9.48}
\end{align}
Der aktuelle Wert des Ausgangsignals ergibt sich bei dem Forward-Euler-Verfahren aus dem letzten Integrationswert \(y[k - 1]\) und der Approximation der Fläche in dem neuen Segment über ein Rechteck der Breite \(T_{A}\) und der Höhe \(u[k - 1]\). Die Transformation der Gleichung in den z-Bereich führt zu der Übertragungsfunktion
\begin{align}
G(z) &= \frac{Y(z)}{U(z)} = \frac{z^{- 1}  T_{A}}{1 - z^{- 1}} = \frac{T_{A}}{z - 1} \label{eq:9.49}
\end{align}
Wieder wird die Übertragungsfunktion \(G(z)\) mit der des Integrators im Laplace-Bereich (Gleichung \eqref{eq:9.38}) verglichen. Die Variable s des Laplace-Bereiches entspricht
\begin{align}
s &= \frac{z - 1}{T_{A}} \label{eq:9.50}
\end{align}
beziehungsweise
\begin{align}
z &= 1 + s  T_{A} \label{eq:9.51}
\end{align}
\begin{example}{Entwurf eines Tiefpasses erster Ordnung nach dem Forward-Euler-Verfahren}
\label{beispiel-entwurf-eines-tiefpasses-erster-ordnung-nach-dem-forward-euler-verfahren}

Für ein PT1-Glied ergibt sich mit dieser Transformation die Übertragungsfunktion
\begin{align}
G(z) &= \frac{Y(z)}{U(z)} = \left. \ G(s) \right|_{s = \frac{z - 1}{T_{A}}} = \frac{1}{1 + T  \frac{z - 1}{T_{A}}} = \frac{T_{A}}{T_{A} - T + T  z} = \frac{T_{A}  z^{- 1}}{\left( T_{A} - T \right)  z^{- 1} + T} \label{eq:9.52}
\end{align}
Das System könnte mit der Differenzengleichung
\begin{align}
y[ k] &= \frac{T_{A}  u[ k - 1] + \left( T - T_{A} \right)  y[ k - 1]}{T} \label{eq:9.53}
\end{align}
implementiert werden.

\end{example}

Auch für das Forward-Euler-Verfahren wird die Abbildung der imaginären Achse der s-Ebene auf die z-Ebene untersucht.
\begin{align}
z &= 1 + j  \omega  T_{A} = 1 + j  \Omega \label{eq:9.54}
\end{align}
Bild \ref{fig:9.11} stellt die Abbildung der imaginären Achse der s Ebene auf die z-Ebene bei Verwendung des Forward-Euler-Verfahrens dar.

\begin{figure}
  \centering
  \includegraphics[alt={Abbildung der imaginären Achse der s-Ebene auf die z-Ebene\\
mit dem Forward-Euler-Verfahren}]{new media/10_11_abbildung_der_imaginren_achse.png}
  \caption{Abbildung der imaginären Achse der s-Ebene auf die z-Ebene\\
mit dem Forward-Euler-Verfahren}
  \label{fig:9.11}
\end{figure}

Nach der Definitionsgleichung der z-Transformation wird die imaginäre Achse der Laplace-Ebene auf den Einheitskreis der z-Ebene abgebildet. Mit dem Forward-Euler-Verfahren wird die imaginäre Achse der Laplace-Ebene im Gegensatz dazu auf die Gerade 1 + jΩ abgebildet. Die negative s-Halbebene wird in den Bereich Re(z) \textless{} 1 abgebildet. Bei dem Forward-Euler-Verfahren kann ein stabiles zeitkontinuierliches System \(G(s)\) in ein instabiles zeitdiskretes System \(G(z)\) abgebildet werden. Aus diesem Grund wird dieses Verfahren im Folgenden nicht weiterverfolgt.

\subsection{Trapezregel und bilineare Transformation}\label{sec:9.2.3}

Bei der Trapezregel wird eine abschnittsweise lineare Funktion zur Approximation der Fläche unter der Kurve \(x(t)\) verwendet.
\begin{align}
y[ k] &= y[ k - 1] + \frac{u[ k] + u[ k - 1]}{2}  T_{A} \label{eq:9.55}
\end{align}
Der aktuelle Wert des Ausgangsignals ergibt sich bei der Trapezregel aus dem letzten Integrationswert \(y[k - 1]\) über ein Trapez der Breite \(T_{A}\) und den beiden Punkten \(u[k - 1]\) und \(u[k]\). Die Transformation der Gleichung in den z-Bereich führt zu der Übertragungsfunktion
\begin{align}
G(z) &= \frac{Y(z)}{U(z)} = \frac{T_{A}}{2}  \frac{1 + z^{- 1}}{1 - z^{- 1}} = \frac{T_{A}}{2}  \frac{z + 1}{z - 1} \label{eq:9.56}
\end{align}
Wieder wird die Übertragungsfunktion \(G(z)\) mit der des Integrators im Laplace-Bereich \eqref{eq:9.38} verglichen. Die Variable s des Laplace-Bereiches entspricht
\begin{align}
s &= \frac{2}{T_{A}}  \frac{z - 1}{z + 1} \label{eq:9.57}
\end{align}
beziehungsweise
\begin{align}
z &= \frac{2 + T_{A}  s}{2 - T_{A}  s} \label{eq:9.58}
\end{align}
Es ergibt sich ein Verfahren zum Entwurf eines zeitdiskreten Systems mit der bilinearen Transformation, das in Tabelle \ref{tab:9.4} zusammengefasst ist.

\begin{table}
  \centering
  \caption{Vorgehen zum Entwurf zeitdiskreter Systeme mit der bilinearen Transformation ohne Prewarping}
  \label{tab:9.4}
  \begin{tabular}{ll}
    \toprule
    Schritt & Beschreibung \\
    \midrule
    1 & Entwurf eines analogen Systems mit der Übertragungsfunktion \(G(s)\) \\
    2 & Transformation des zeitkontinuierlichen Systems \(G(s)\) in ein zeitdiskretes System mit der bilinearen Transformation, Durchführung der Substitution \\
    & \(s = \frac{2}{T_{A}}  \frac{z - 1}{z + 1}\) \\
    3 & Aufstellen der Differenzengleichung durch Rücktransformation der Gleichung \(G(z)\) in den Zeitbereich \\
    & \(G(z) = \frac{Y(z)}{U(z)}\) \\
    4 & Auflösen der Differenzengleichung nach \(y[k]\) \\
    \bottomrule
  \end{tabular}
\end{table}

\begin{example}{Entwurf eines Tiefpasses erster Ordnung mit der bilinearen Transformation}
\label{beispiel-entwurf-eines-tiefpasses-erster-ordnung-mit-der-bilinearen-transformation}

Für ein PT1-Glied ergibt sich mit der bilinearen Transformation die Übertragungsfunktion
\begin{align}
{G(z) = \frac{Y(z)}{U(z)} = \left. \ G(s) \right|_{s = \frac{2}{T_{A}}  \frac{z - 1}{z + 1}} = \frac{1}{1 + T  \frac{2}{T_{A}}  \frac{z - 1}{z + 1}} = \frac{z  T_{A} + T_{A}}{z  T_{A} + T_{A} + 2  T  z - 2  T}
}{= \frac{z  T_{A} + T_{A}}{z  \left( T_{A} + 2  T \right) + T_{A} - 2  T} = \frac{T_{A} + T_{A}  z^{- 1}}{\left( T_{A} + 2  T \right) + \left( T_{A} - 2  T \right)z^{- 1}}} \label{eq:9.59}
\end{align}
Das System könnte mit der Differenzengleichung
\begin{align}
y[ k] &= \frac{T_{A}  u[ k] + T_{A}  u[ k - 1] - \left( T_{A} - 2  T \right)  y[ k - 1]}{T_{A} + 2  T} \label{eq:9.60}
\end{align}
implementiert werden.

\end{example}

Durch Einsetzen der Bedingung s = jω wird wieder die Abbildung der imaginären Achse der s-Ebene auf die z-Ebene untersucht.
\begin{align}
z &= \frac{2 + T_{A}  j  \omega}{2 - T_{A}  j  \omega} = e^{j  2  \arctan\left( \frac{\omega  T_{A}}{2} \right)} = e^{j  2  \arctan\left( \frac{\Omega}{2} \right)} \label{eq:9.61}
\end{align}
Bild \ref{fig:9.12} stellt die Abbildung der imaginären Achse der s-Ebene auf die z-Ebene bei Approximation der Variable z über eine algebraische Gleichung in s dar.

\begin{figure}
  \centering
  \includegraphics[alt={Abbildung der imaginären Achse der s-Ebene auf die z-Ebene\\
bei Approximation der Variable z über eine algebraische Gleichung in s}]{new media/10_12_abbildung_der_imaginren_achse.png}
  \caption{Abbildung der imaginären Achse der s-Ebene auf die z-Ebene\\
bei Approximation der Variable z über eine algebraische Gleichung in s}
  \label{fig:9.12}
\end{figure}

Die bilineare Transformation bildet die imaginäre Achse der s-Ebene auf den Einheitskreis der\\
z-Ebene ab, die linke s-Halbebene wird in das Innere des Einheitskreises der z-Ebene abgebildet. Damit bleibt ein stabiles System bei der bilinearen Transformation ein stabiles System. Ein Vergleich mit dem bei der Herleitung gebrauchten Ausdruck
\begin{align}
z &= e^{j  \Omega} \label{eq:9.62}
\end{align}
führt zu dem Zusammenhang zwischen der Kreisfrequenz ω des zeitkontinuierlichen Systems und der normierten Kreisfrequenz \(Ω_{B}\) des zeitdiskreten Systems, das sich aus der bilinearen Transformation ergibt.
\begin{align}
\Omega_{B} &= 2  \arctan\left( \frac{\omega  T_{A}}{2} \right) = 2  \arctan\left( \frac{\Omega}{2} \right) \label{eq:9.63}
\end{align}
beziehungsweise
\begin{align}
\omega &= \frac{2}{T_{A}}  \tan\left( \frac{\Omega_{B}}{2} \right) \label{eq:9.64}
\end{align}
oder
\begin{align}
\Omega &= 2  \tan\left( \frac{\Omega_{B}}{2} \right) \label{eq:9.65}
\end{align}
Die Frequenz wird durch die bilineare Transformation nicht linear, sondern verzerrt abgebildet. In Bild \ref{fig:9.13} wird diese Verzerrung der Frequenzachse grafisch dargestellt.

\begin{figure}
  \centering
  \includegraphics[alt={Visualisierung der Verzerrung der Frequenzachse durch die bilineare Transformation}]{new media/10_13_visualisierung_der_verzerrung.png}
  \caption{Visualisierung der Verzerrung der Frequenzachse durch die bilineare Transformation}
  \label{fig:9.13}
\end{figure}

Die Visualisierung der Verzerrung ist im Bereich kleiner Frequenzen gering und steigt mit steigender Frequenz an. Sie ist mathematisch eindeutig und kann beim Systementwurf berücksichtigt werden. Dazu erfolgt der Systementwurf mit einer Spezifikation, die die bei der Diskretisierung stattfindende Frequenzverschiebung berücksichtigt. Das Verfahren wird in Abschnitt \ref{sec:9.2.4} beschrieben.

\begin{example}{Zeitdiskrete Beschreibung des Aufheizverhaltens eines Transistors}
\label{beispiel-zeitdiskrete-beschreibung-des-aufheizverhaltens-eines-transistors}

Das Aufheizverhalten eines Transistors kann über die Differentialgleichung
\begin{align}
C_{TH}  \frac{dT}{dt} &= p_{EL}(t) - \frac{1}{R_{TH}}  T(t) \label{eq:9.66}
\end{align}
beschrieben werden. Dabei bezeichnet T die Temperaturdifferenz zur Umgebungstemperatur gemessen in K, \(C_{TH}\) ist die thermische Wärmekapazität, \(R_{TH}\) der thermische Widerstand und \(p_{EL}\) die elektrische Leistung in W, die den Transistor aufheizt. In der Anordnung beträgt die thermische Kapazität\\
\(C_{TH} = 0.2\) Ws/K und der thermische Widerstand \(R_{TH} = 5\) K/W.

Das Aufheizverhalten soll in Abhängigkeit der elektrischen Leistung \(p_{EL}(t)\) zeitdiskret simuliert werden. Zum Zeitpunkt \(t = 0\) hat das System die Temperatur der Umgebung T(0) = 0. Für den Zeitraum\\
t \textless{} 0 ist die elektrische Leistung \(p_{EL}(t)\) = 0, für t \textgreater{} 0 ist die elektrische Leistung zunächst eine beliebige Funktion der Zeit.

Transformation der Differentialgleichung in den Laplace-Bereich führt unter Berücksichtigung der Differentiationsregel und wegen der verschwindenden Anfangsbedingungen zu
\begin{align}
\left( R_{TH}  C_{TH}  s + 1 \right)  T(s) &= R_{TH}  P_{EL}(s) \label{eq:9.67}
\end{align}
Aus der Gleichung im Laplace-Bereich ergibt sich die Übertragungsfunktion
\begin{align}
\frac{T(s)}{P_{EL}(s)} &= \frac{R_{TH}}{R_{TH}  C_{TH}  s + 1} \label{eq:9.68}
\end{align}
Mit der bilinearen Transformation ergibt sich aus der Übertragungsfunktion im Laplace-Bereich die Übertragungsfunktion im z-Bereich zu
\begin{align}
{\frac{T\left(z\right)}{P_{EL} \left(z\right)} } &=  {\left. \frac{R_{TH} }{R_{TH}  C_{TH}  s+1} \right|_{s=\frac{2}{T_{A} }  \frac{z-1}{z+1} } =\frac{R_{TH} }{R_{TH}  C_{TH}  \frac{2}{T_{A} }  \frac{z-1}{z+1} +1} } \\
{} &=  {\frac{R_{TH}  T_{A}  \left(z+1\right)}{R_{TH}  C_{TH}  2 \left(z-1\right)+T_{A}  \left(z+1\right)} } \\
{} &=  {\frac{R_{TH}  T_{A}  z+R_{TH}  T_{A} }{\left(R_{TH}  C_{TH}  2+T_{A} \right) z+\left(T_{A} -R_{TH}  C_{TH}  2\right)} } \\
{} &=  {\frac{R_{TH}  T_{A} +R_{TH}  T_{A}  z^{-1} }{\left(R_{TH}  C_{TH}  2+T_{A} \right) z+\left(T_{A} -R_{TH}  C_{TH}  2\right) z^{-1} } }
 \label{eq:9.69}
\end{align}
Nach dem Ausmultiplizieren
\begin{align}
{T(z)  \left( R_{TH}  C_{TH}  2 + T_{A} \right) + T(z)  \left( T_{A} - R_{TH}  C_{TH}  2 \right)  z^{- 1}
}{= R_{TH}  T_{A}  P_{EL}(z) + R_{TH}  T_{A}  z^{- 1}  P_{EL}(z)} \label{eq:9.70}
\end{align}
kann das System als Differenzengleichung dargestellt werden
\begin{align}
{T[ k]  \left( R_{TH}  C_{TH}  2 + T_{A} \right) + T[ k - 1]  \left( T_{A} - R_{TH}  C_{TH}  2 \right)
}{= R_{TH}  T_{A}  p_{EL}[ k] + R_{TH}  T_{A}  p_{EL}[ k - 1]} \label{eq:9.71}
\end{align}
Auflösen nach \(T[k]\) ergibt
\begin{align}
T[ k] &= \frac{R_{th}  T_{A}  p_{el}[ k] + R_{th}  T_{A}  p_{el}[ k - 1] - T[ k - 1]  \left( T_{A} - R_{th}  C_{th}  2 \right)}{R_{th}  C_{th}  2 + T_{A}} \label{eq:9.72}
\end{align}
Auf Basis der - \SI{20}{\decibel}-Grenzfrequenz des analogen Systems wird die Abtastzeit \(T_{A}\) bestimmt. Das zeitkontinuierliche System hat einen einfachen reellen Pol in der negativen Halbebene, es ist demnach stabil. Deshalb ergibt sich der Frequenzgang zu
\begin{align}
G(\omega) &= \left. \ \frac{T(s)}{P_{EL}(s)} \right|_{s = j  \omega} = \frac{R_{TH}}{R_{TH}  C_{TH}  j  \omega + 1} \label{eq:9.73}
\end{align}
Um die - \SI{20}{\decibel} Grenzfrequenz ausrechnen zu können, wird zunächst der Amplitudengang an der Stelle \(ω = 0\) bestimmt.
\begin{align}
A(0) &= R_{TH} \label{eq:9.74}
\end{align}
Daraus ergibt sich die Grenzfrequenz aus der Bedingung
\begin{align}
A\left( \omega_{G} \right) &= \left| \frac{R_{TH}}{R_{TH}  C_{TH}  j  \omega_{G} + 1} \right| = \frac{A(0)}{10} = \frac{R_{TH}}{10} \label{eq:9.75}
\end{align}
Daraus ergibt sich die Bedingung
\begin{align}
100 &= R_{TH}^{2}  C_{TH}^{2}  \omega_{G}^{2} + 1 \label{eq:9.76}
\end{align}
Die Grenzfrequenz muss positiv sein, damit ergibt sich mit den Zahlenwerten aus der Aufgabenstellung
\begin{align}
\omega_{G} &= \frac{\sqrt{100 - 1}}{R_{th}  C_{th}} = \frac{\sqrt{99}}{\SI{1}{\second}} = 9.9499\ \frac{rad}{s} \approx 10\ \frac{rad}{s} \label{eq:9.77}
\end{align}
Mit der berechneten Grenzfrequenz ergibt sich für die Abtastfrequenz
\begin{align}
\omega_{A} &= 2  \omega_{G} = 20\ \frac{rad}{s} = \frac{2  \pi}{T_{A}} \label{eq:9.78}
\end{align}
und damit für die Abtastzeit
\begin{align}
T_{A} &= \frac{2  \pi}{20}\ s = \frac{\pi}{10}\ s \label{eq:9.79}
\end{align}
Zur Berechnung des Ausgangssignals \(T(k)\) muss das Eingangssignal diskretisiert und die Differenzengleichung unter Berücksichtigung der Anfangsbedingung rekursiv berechnet werden. Es ergibt sich folgendes Einschwingverhalten.

\begin{fixedfigure}
  \includegraphics[alt={Vergleich des simulierten Aufheizverhaltens für unterschiedliche Abtastzeiten\\
\(T_{A1} = π\)/10 und \(T_{A2} = π\)/100}]{new media/10_14_vergleich_des_simulierten_aufh.png}
  \captionof{figure}{Vergleich des simulierten Aufheizverhaltens für unterschiedliche Abtastzeiten\\
\(T_{A1} = π\)/10 und \(T_{A2} = π\)/100}
\end{fixedfigure}

In der Grafik ist zum Vergleich das Signal dargestellt, dass sich bei einer 10-mal kleineren Abtastzeit ergibt. Mit sinkender Abtastzeit steigt die Genauigkeit der Simulation.

\end{example}

\subsection{Bilineare Transformation mit Prewarping}\label{sec:9.2.4}

Die in Abschnitt \ref{sec:9.2.3} beschriebene bilineare Transformation kann so modifiziert werden, dass sie die Frequenzverschiebung an einer Frequenz \(ω_{0}\) kompensiert. Sie ergibt sich aus der Substitution
\begin{align}
s &= V  \frac{z - 1}{z + 1} \label{eq:9.80}
\end{align}
Dabei ist V ein noch zu bestimmender Parameter. Die Frequenzverzerrung ergibt sich nach einer Transformation in den Frequenzbereich mit s = jω und z = \(e^{jΩp}\) zu
\begin{align}
{j \omega } &=  {V \frac{e^{j \Omega _{P} } -1}{e^{j \Omega _{P} } +1} =V \frac{e^{j \frac{\Omega _{P} }{2} } }{e^{j \frac{\Omega _{P} }{2} } }  \frac{e^{j \frac{\Omega _{P} }{2} } -e^{-j \frac{\Omega _{P} }{2} } }{e^{j \frac{\Omega _{P} }{2} } +e^{j- \frac{\Omega _{P} }{2} } } =V \frac{2 j}{2}  \frac{\sin \left(\frac{\Omega _{P} }{2} \right)}{\cos \left(\frac{\Omega _{P} }{2} \right)} } \\
{} &=  {j V tan\left(\frac{\Omega _{P} }{2} \right)=j V tan\left(\frac{\omega _{P}  T_{A} }{2} \right)}
 \label{eq:9.81}
\end{align}
Damit der Frequenzgang der zeitkontinuierlichen und zeitdiskreten Implementierung an der Frequenz \(ω_{0}\) übereinstimmen, muss damit die Bedingung
\begin{align}
\omega_{0} &= V  tan\left( \frac{\omega_{0}  T_{A}}{2} \right) \label{eq:9.82}
\end{align}
beziehungsweise
\begin{align}
V &= \frac{\omega_{0}}{tan\left( \frac{\omega_{0}  T_{A}}{2} \right)} \label{eq:9.83}
\end{align}
erfüllt werden. Es ergibt sich die bilineare Transformation mit Prewarping:
\begin{align}
s &= V  \frac{z - 1}{z + 1} = \frac{\omega_{0}}{tan\left( \frac{\omega_{0}  T_{A}}{2} \right)}  \frac{z - 1}{z + 1} \label{eq:9.84}
\end{align}
Auch die bilineare Transformation mit Prewarping bildet die imaginäre Achse der s-Ebene auf den Einheitskreis der z-Ebene ab, die linke s-Halbebene wird in das Innere des Einheitskreises der z-Ebene abgebildet. Damit bleibt ein stabiles System bei der bilinearen Transformation mit Prewarping ein stabiles System. Gleichung \eqref{eq:9.81} wird dazu verwendet, den Zusammenhang zwischen der Kreisfrequenz ω des zeitkontinuierlichen Systems und der normierten Kreisfrequenz \(Ω_{P}\) des zeitdiskreten Systems mit Prewarping zu bestimmen.
\begin{align}
\omega &= V  tan\left( \frac{\Omega_{P}}{2} \right) \label{eq:9.85}
\end{align}
Auflösen nach der normierten Kreisfrequenz \(Ω_{P}\) führt zu
\begin{align}
\Omega_{P} &= 2  \arctan\left( \frac{\omega}{V} \right) = 2  \arctan\left( tan\left( \frac{\omega_{0}  T_{A}}{2} \right)  \frac{\omega}{\omega_{0}} \right) \label{eq:9.86}
\end{align}
An der Frequenz \(ω = ω_{0}\) vereinfacht sich der Ausdruck zu
\begin{align}
\Omega_{P0} &= 2  \arctan\left( \frac{\omega_{0}}{V} \right) = 2  \arctan\left( tan\left( \frac{\omega_{0}  T_{A}}{2} \right)  \frac{\omega_{0}}{\omega_{0}} \right) = 2  \frac{\omega_{0}  T_{A}}{2} = \omega_{0}  T_{A} = \Omega_{0} \label{eq:9.87}
\end{align}
und die beiden normierten Kreisfrequenzen \(Ω_{P0}\) und \(Ω_{0}\) sind identisch. In Bild \ref{fig:9.15} wird die Verzerrung der Frequenzachse bei der bilinearen Transformation mit Prewarping für \(Ω_{0} = π\)/2 grafisch dargestellt.

\begin{figure}
  \centering
  \includegraphics[alt={Verzerrung der Frequenzachse durch die bilineare Transformation mit Prewarping für \(Ω_{0} = π\)/2}]{new media/10_15_verzerrung_der_frequenzachse_d.png}
  \caption{Verzerrung der Frequenzachse durch die bilineare Transformation mit Prewarping für \(Ω_{0} = π\)/2}
  \label{fig:9.15}
\end{figure}

Die nichtlineare Frequenzverzerrung wird mit Prewarping nicht generell aufgehoben, aber an der Stelle \(ω_{0}\) kompensiert.

\begin{example}{Entwurf eines Tiefpasses erster Ordnung mit der bilinearen Transformation und Prewarping}
\label{beispiel-entwurf-eines-tiefpasses-erster-ordnung-mit-der-bilinearen-transformation-und-prewarping}

Für ein PT1-Glied ergibt sich mit der bilinearen Transformation mit Prewarping die Übertragungsfunktion
\begin{align}
{G(z) = \frac{Y(z)}{U(z)} = \left. \ G(s) \right|_{s = V  \frac{z - 1}{z + 1}} = \frac{1}{1 + T  V  \frac{z - 1}{z + 1}} = \frac{z + 1}{z + 1 + T  V  z - T  V}
}{= \frac{z + 1}{z  (1 + T  V) + 1 - T  V}} \label{eq:9.88}
\end{align}
Das System könnte mit der Differenzengleichung
\begin{align}
y[ k] &= \frac{u[ k] + u[ k - 1] - (1 - T  V)  y[ k - 1]}{1 + T  V} \label{eq:9.89}
\end{align}
implementiert werden. Der Faktor V wird so gewählt, dass die 3-dB-Grenzfrequenzen des zeitkontinuierlichen und des zeitdiskreten Filters identisch sind. Mit \(ω_{G} = 1\)/\SI{2}{\radian\per\second} und \(T_{A} = \SI{1}{\second}\) ergibt sich für V der Wert
\begin{align}
V &= \frac{\omega_{G}}{tan\left( \frac{\omega_{G}  T_{A}}{2} \right)} = \frac{1}{2  tan\left( \frac{1}{4} \right)} = \text{1.9582} \label{eq:9.90}
\end{align}
Damit lautet die Übertragungsfunktion für das mit Prewarping entworfene zeitdiskrete Filter
\begin{align}
G(z) &= \frac{z + 1}{z  (1 + V  T) + 1 - V  T} = \frac{z + 1}{\text{4.9163}  z - \text{2.9163}} \label{eq:9.91}
\end{align}
Das Ergebnis kann mit der Berechnung der Amplitudengänge verifiziert werden.

\begin{fixedfigure}
  \includegraphics[alt={Amplitudengang des zeitkontinuierlichen Tiefpasses und der über bilineare Transformation\\
ohne und mit Prewarping entworfenen Systeme}]{new media/10_16_amplitudengang_des_zeitkontinu.png}
  \captionof{figure}{Amplitudengang des zeitkontinuierlichen Tiefpasses und der über bilineare Transformation\\
ohne und mit Prewarping entworfenen Systeme}
\end{fixedfigure}

An der Grenzfrequenz \(ω_{G} = 1\)/\SI{2}{\radian\per\second} stimmen der Amplitudengang des zeitkontinuierlichen Systems und des über Prewarping entworfenen Systems exakt überein. Im Gegensatz dazu weist das über bilineare Transformation ohne Prewarping entworfene System eine Frequenzverschiebung auf.

\end{example}

Damit ergibt sich ein Verfahren zum Entwurf eines zeitdiskreten Systems mit der bilinearen Transformation und Prewarping, das in Tabelle \ref{tab:9.5} zusammengefasst ist.

\begin{table}
  \centering
  \caption{Vorgehen zum Entwurf rekursiver Filter mit der bilinearen Transformation und Prewarping}
  \label{tab:9.5}
  \begin{tabular}{ll}
    \toprule
    Schritt & Beschreibung \\
    \midrule
    1 & Entwurf eines zeitkontinuierlichen Filters \(G(s)\) \\
    2 & Bestimmung der Konstanten V über die Frequenz \(ω_{0}\), bei der das zeitkontinuierliche und zeitdiskrete System denselben Wertz für den Frequenzgang aufweisen sollen \\
    & \(V = \frac{\omega_{0}}{tan\left( \frac{\omega_{0}  T_{A}}{2} \right)}\) \\
    3 & Transformation des zeitkontinuierlichen Filters \(G(s)\) in ein zeitdiskretes Filter mit der bilinearen Transformation, Durchführung der Substitution \\
    & \(s = V  \frac{z - 1}{z + 1}\) \\
    3 & Aufstellen der Differenzengleichung durch Rücktransformation der Gleichung \(G(z)\) in den Zeitbereich \\
    & \(G(z) = \frac{Y(z)}{U(z)}\) \\
    4 & Auflösen der Differenzengleichung nach \(y[k]\) \\
    \bottomrule
  \end{tabular}
\end{table}

Die bilineare Transformation mit Prewarping ist das zweite Verfahren, das in der Regelungstechnik für die zeitdiskrete Nachbildung analoger Systeme verwendet wird. Da die Genauigkeit der Transformation in regelungstechnischen Anwendungen an der Durchtrittsfrequenz wichtig ist, wird ein Prewarping mit dieser Frequenz durchgeführt. Auch bei der Berechnung zeitdiskreter Filter wird die bilineare Transformation mit Prewarping eingesetzt. Dabei wird das Prewarping mit der spezifizierten Grenzfrequenz durchgeführt.

Der Vorteil der bilinearen Transformation liegt in der einfachen Berechnung über die Substitution
\begin{align}
s &= V  \frac{z - 1}{z + 1} \label{eq:9.92}
\end{align}
Der Nachteil ergibt sich aus der groben Approximation des Integrals mit einer Trapezregel.

\subsection{Entwurf eines IIR-Filters mit der bilinearen Transformation}\label{sec:entwurf-eines-iir-filters-mit-der-bilinearen-transformation}

Zur Verdeutlichung des Vorgehens wird das Beispiel des Butterworth-Tiefpasses aus Abschnitt \ref{sec:9.1.4} aufgegriffen. Zur Kompensation der Frequenzverschiebung wird die Grenzfrequenz gemäß Gleichung \hyperref[_Ref382407345]{\textbf{Fehler! Verweisquelle konnte nicht gefunden werden.}} umgerechnet zu
\begin{align}
\omega_{GP} &= \frac{2}{T_{A}}  \tan\left( \frac{\omega_{G}  T_{A}}{2} \right) = \SI{200}{\kilo\hertz}  \tan\left( \frac{2  \pi  10^{4}  10^{- 5}}{2} \right) = \SI{64.98}{\kilo\hertz} \label{eq:9.93}
\end{align}
Das zeitkontinuierliche Filter hat die Übertragungsfunktion
\begin{align}
{G\left(s\right)} &=  {\frac{\omega _{GP}^{3} }{\left(s+\omega _{GP} \right) \left(s+\left(\frac{1}{2} -j \frac{\sqrt{3} }{2} \right) \omega _{GP} \right) \left(s+\left(\frac{1}{2} +j \frac{\sqrt{3} }{2} \right) \omega _{GP} \right)} } \\
{} &=  {\frac{\omega _{GP}^{3} }{s^{3} +\SI{2}{\second}^{2}  \omega _{GP} +\SI{2}{\second} \omega _{GP}^{2} +\omega _{GP}^{3} } =\frac{1}{\left(\frac{s}{\omega _{GP} } \right)^{3} +2 \left(\frac{s}{\omega _{GP} } \right)^{2} +2 \left(\frac{s}{\omega _{GP} } \right)+1} }
 \label{eq:9.94}
\end{align}
Mit der bilinearen Transformation wird das zeitkontinuierliche System in ein zeitdiskretes Filter mit der Übertragungsfunktion
\begin{align}
{G\left(z\right)} &=  {\left. \frac{1}{\left(\frac{s}{\omega _{GP} } \right)^{3} +2 \left(\frac{s}{\omega _{GP} } \right)^{2} +2 \left(\frac{s}{\omega _{GP} } \right)+1} \right|_{s=\frac{2}{T_{A} }  \frac{z-1}{z+1} } } \\
{} &=  {{0.0181} \frac{\left(z+1\right)^{3} }{z^{3} -1.760 z^{2} +1.183 z-0.2781} }
 \label{eq:9.95}
\end{align}
überführt. Bild \ref{fig:9.17} stellt das Pol-Nullstellendiagramm für das über die bilineare Transformation entworfene Butterworth-Filter dar.

\begin{figure}
  \centering
  \includegraphics[alt={Pol-Nullstellendiagramm für das über die bilineare Transformation\\
entworfene Butterworth-Filter dritter Ordnung}]{new media/10_17_polnullstellendiagramm_fr_das.png}
  \caption{Pol-Nullstellendiagramm für das über die bilineare Transformation\\
entworfene Butterworth-Filter dritter Ordnung}
  \label{fig:9.17}
\end{figure}

Während die Pole eine ähnliche Lage haben wie bei dem impulsinvarianten Entwurf, liegt an dem Punkt z = - 1 eine Nullstelle dritter Ordnung. Dadurch weist der Amplitudengang an der Stelle Ω = ± π exakt den Wert 0 auf. Die Pole der Übertragungsfunktion liegen auch bei diesem Verfahren im Inneren des Einheitskreises, das System ist demnach stabil. Die Verstärkung des zeitdiskreten Systems errechnet sich zu
\begin{align}
G(1) &= \text{0.0181}  \frac{2^{3}}{1 - \text{1.760} + \text{1.183} - \text{0.2781}} = 1 \label{eq:9.96}
\end{align}
Der Frequenzgang \(G(Ω)\) ergibt sich durch die Substitution z = \(e^{jΩ}\) zu
\begin{align}
G(\Omega) &= \text{0.0181}  \frac{e^{j  3  \Omega} + 3  e^{j  2  \Omega} + 3  e^{j  \Omega} + 1}{e^{j  3  \Omega} - \text{1.760}  e^{j  2  \Omega} + \text{1.183}  e^{j  \Omega} - \text{0.2781}} \label{eq:9.97}
\end{align}
Bild \ref{fig:9.18} zeigt den Frequenzgang des zeitdiskreten Systems im Vergleich zu dem Frequenzgang des zeitkontinuierlichen Systems im linearen Maßstab.

\begin{figure}
  \centering
  \includegraphics[alt={Frequenzgang des über die bilineare Transformation entworfenen\\
Butterworth-Filters dritter Ordnung}]{new media/10_18_frequenzgang_des_ber_die_bilin.png}
  \caption{Frequenzgang des über die bilineare Transformation entworfenen\\
Butterworth-Filters dritter Ordnung}
  \label{fig:9.18}
\end{figure}

Der Vergleich des zeitkontinuierlichen Filters mit dem zeitdiskreten Filter zeigt erwartungsgemäß die die Frequenzverschiebung gemäß der Beziehung
\begin{align}
\Omega_{B} &= 2  \arctan\left( \frac{\omega  T_{A}}{2} \right) \label{eq:9.98}
\end{align}
Durch Prewarping sind die \SI{3}{\decibel}-Grenzfrequenzen des zeitkontinuierlichen und des zeitdiskreten Filters identisch.

Mit der Übertragungsfunktion \(G(z)\) kann das Filter als Differenzengleichung realisiert werden. Ausgehend von der Übertragungsfunktion ergibt sich nach der Division von Zähler und Nenner mit der höchsten Potenz von z die Gleichung
\begin{align}
G(z) &= \frac{Y(z)}{U(z)} = \text{0.0167}  \frac{1 + 3  z^{- 1} + 3  z^{- 2} + z^{- 3}}{1 - 1.798  z^{- 1} + 1.221  z^{- 2} - 0.2898  z^{- 3}} \label{eq:9.99}
\end{align}
und damit die Differenzengleichung
\begin{align}
{y\left[k\right]} &=  {1.798 y\left[k-1\right]-1.221 y\left[k-2\right]+0.2898 y\left[k-3\right]} \\
{} & {} & {+0.0167 u\left[k\right]+0.0501 u\left[k-1\right]+0.0501 u\left[k-2\right]+0.0167 u\left[k-3\right]}
 \label{eq:9.100}
\end{align}
Bild \ref{fig:9.19} stellt die Impuls- und Sprungantwort des zeitdiskreten Systems im Vergleich zur Impuls- und Sprungantwort des zeitkontinuierlichen Systems. Wieder werden die zeitdiskreten Signale entsprechend normiert.

\begin{figure}
  \centering
  \includegraphics[alt={Impuls- und Sprungantwort für das über die bilineare Transformation\\
entworfenen Butterworth-Filter dritter Ordnung}]{new media/10_19_impuls_und_sprungantwort_fr_da.png}
  \caption{Impuls- und Sprungantwort für das über die bilineare Transformation\\
entworfenen Butterworth-Filter dritter Ordnung}
  \label{fig:9.19}
\end{figure}

Im Gegensatz zum impulsinvarianten Entwurf stimmen zeitkontinuierliche und zeitdiskrete Impulsantwort nicht mehr exakt überein. Auch die Sprungantworten unterscheiden sich.

\section{Analytischer Entwurf von FIR-Filtern}\label{sec:9.3}

Nichtrekursive Filter besitzen unter den in Abschnitt \ref{sec:8.4.6} beschriebenen Voraussetzungen eine lineare Phase. Die Vorteile, die sich daraus ergeben, sind in Abschnitt \ref{sec:8.4.7} dargestellt. In diesem Abschnitt wird der Entwurf nicht rekursiver Filter im Zeit- und Frequenzbereich dargestellt. Der Filterentwurf im Zeitbereich erfolgt über ein Abtasten der Impulsantwort. Der Filterentwurf im Frequenzbereich basiert auf der Approximation des gewünschten Frequenzgangs über eine Fourier-Reihe. Beide Verfahren führen zu demselben Ergebnis. Durch das Begrenzen der Impulsantwort auf einen endlichen Zeitbereich entsteht in Kombination mit dem Ziel, einen idealen Tiefpass zu realisieren, ein Überschwingen des Frequenzgangs. Der Effekt ist als Gibbscher Effekt bekannt. Zur Kompensation des Effektes kann eine Gewichtung der Abtastwerte vorgenommen werden, die als Fensterung bezeichnet wird.

\subsection{Idealer Filtercharakteristiken}\label{sec:9.3.1}

Bei der Herleitung des Frequenzgangs in Kapitel \ref{sec:8.1.3} wird gezeigt, dass sich das Spektrum des Ausgangssignals \(Y(Ω)\) aus dem Produkt der Spektren von Impulsantwort \(G(Ω)\) und Eingangssignal \(U(Ω)\) ergibt.
\begin{align}
Y(\Omega) &= G(\Omega)  U(\Omega) \label{eq:9.101}
\end{align}
Der Frequenzgang \(G(Ω)\) kann dazu genutzt werden, Spektralanteile zu unterdrücken oder zu verstärken. Die Verstärkung beziehungsweise die Unterdrückung von Frequenzbereichen kann am Amplitudengang abgelesen werden. Analog zu zeitkontinuierlichen Filtern ergeben sich vier unterschiedliche ideale Filter, nämlich Tiefpass, Hochpass, Bandpass und Bandsperre. Im Gegensatz zu zeitkontinuierlichen Systemen muss dabei auf die Periodizität des Frequenzgangs geachtet werden.

\subsubsection{Idealer Tiefpass}\label{sec:idealer-tiefpass}

Das ideale Tiefpass-Filter weist eine Übertragungsfunktion auf, die niederfrequente Anteile bis zu einer Grenzfrequenz \(Ω_{G}\) ungedämpft durchlässt und hochfrequente Anteile oberhalb der Grenzfrequenz \(Ω_{G}\) sperrt. Die Übertragungsfunktion des idealen Tiefpass-Filters lautet
\begin{align}
G_{TP}(\Omega) = \left\{ \begin{aligned}
 & 1\text{   f}\text{ü}\text{r }|\Omega| < \Omega_{G} \\
 & 0\text{  f}\text{ü}\text{r }\Omega_{G} < |\Omega| < \pi\ 
\end{aligned} \right.\ \label{eq:9.102}
\end{align}
Die Impulsantwort des idealen Tiefpass-Filters lautet
\begin{align}
g_{TP}[ k] &= \frac{\sin\left( \Omega_{G}  k \right)}{\pi  k} \label{eq:9.103}
\end{align}
Amplitudengang und Impulsantwort des idealen Tiefpass-Filters sind in Bild \ref{fig:9.20} dargestellt.

\begin{figure}
  \centering
  \includegraphics[alt={Amplitudengang und Impulsantwort des idealen Tiefpass-Filters mit \(Ω_{G} = π\)/3}]{new media/10_20_amplitudengang_und_impulsantwo.png}
  \caption{Amplitudengang und Impulsantwort des idealen Tiefpass-Filters mit \(Ω_{G} = π\)/3}
  \label{fig:9.20}
\end{figure}

Die Impulsantwort \(g_{TP}\){[}k{]} ist nicht auf den Bereich k ≥ 0 beschränkt, sie ist demnach nicht kausal und besitzt eine unendliche Länge. Das ideale Tiefpass-Filter ist deshalb technisch nicht realisierbar.

Neben der Eigenschaft der frequenzselektiven Dämpfung muss das ideale Tiefpass-Filter zumindest im Durchlassbereich einen Phasengang von ϕ = 0 aufweisen. Es wird sich zeigen, dass kausale Systeme immer eine Phase ϕ(Ω) ≠ 0 aufweisen. Es wird sich außerdem zeigen, dass eine lineare Phase im Durchlassbereich nur zu einer Verschiebung des Signals führt und das Signal darüber hinaus nicht verzerrt wird. Dadurch kann die Forderung an den Phasengang abgeschwächt werden.

\subsubsection{Idealer Hochpass}\label{sec:idealer-hochpass}

Das ideale Hochpass-Filter weist eine Übertragungsfunktion auf, die niederfrequente Anteile bis zu einer Grenzfrequenz \(Ω_{G}\) sperrt und hochfrequente Anteile oberhalb der Grenzfrequenz \(Ω_{G}\) ungedämpft durchlässt. Die Übertragungsfunktion des idealen Hochpass-Filters lautet damit
\begin{align}
G_{HP}(\Omega) = \left\{ \begin{aligned}
 & \text{0   f}\text{ü}\text{r }|\Omega| < \Omega_{G} \\
 & \text{1   f}\text{ü}\text{r }\Omega_{G} < |\Omega| < \pi\ 
\end{aligned} \right.\ \label{eq:9.104}
\end{align}
Die Übertragungsfunktion kann über die Übertragungsfunktion des Tiefpass-Filters ausgedrückt werden.
\begin{align}
G_{HP}(\Omega) &= 1 - G_{TP}(\Omega) \label{eq:9.105}
\end{align}
Die Impulsantwort des idealen Hochpass-Filters ergibt sich nach den Rechenregeln der Fourier-Transformation von Signalfolgen zu
\begin{align}
g_{HP}[ k] &= \delta[ k] - g_{TP}[ k] = \delta[ k] - \frac{\sin\left( \Omega_{G}  k \right)}{\pi  k} \label{eq:9.106}
\end{align}
Amplitudengang und Impulsantwort des idealen Hochpass-Filters sind in Bild \ref{fig:9.21} dargestellt.

\begin{figure}
  \centering
  \includegraphics[alt={Amplitudengang und Impulsantwort des idealen Hochpass-Filters mit \(Ω_{G} = π\)/3}]{new media/10_21_amplitudengang_und_impulsantwo.png}
  \caption{Amplitudengang und Impulsantwort des idealen Hochpass-Filters mit \(Ω_{G} = π\)/3}
  \label{fig:9.21}
\end{figure}

Auch die Impulsantwort \(g_{HP}\){[}k{]} ist nicht auf den Bereich k ≥ 0 beschränkt und deshalb technisch nicht realisierbar.

\subsubsection{Ideale Bandsperre und idealer Bandpass}\label{sec:ideale-bandsperre-und-idealer-bandpass}

Neben dem idealen Hoch- und Tiefpass existieren die ideale Bandsperre und der ideale Bandpass. Die ideale Bandsperre sperrt in einem Bereich zwischen den Frequenzen \(Ω_{G1}\) und \(Ω_{G2}\). Sie besitzt die Übertragungsfunktion
\begin{align}
G_{BS}(\Omega) = \left\{ \begin{aligned}
 & \text{1   f}\text{ü}\text{r }|\Omega| < \Omega_{G1}\text{ und f}\text{ü}\text{r }\Omega_{G2} < |\Omega| < \pi \\
 & \text{0   f}\text{ü}\text{r }\Omega_{G1} < |\Omega| < \Omega_{G2}\ 
\end{aligned} \right.\ \label{eq:9.107}
\end{align}
Der ideale Bandpass lässt Signalanteile in einem Bereich zwischen den Frequenzen \(Ω_{G1}\) und \(Ω_{G2}\) durch. Er hat folgende Übertragungsfunktion:
\begin{align}
G_{BP}(\Omega) = \left\{ \begin{aligned}
 & \text{1   f}\text{ü}\text{r }\Omega_{G1} < |\Omega| < \Omega_{G2} \\
 & \text{0   f}\text{ü}\text{r }|\Omega| < \Omega_{G1}\text{ und f}\text{ü}\text{r }\Omega_{G2} < |\Omega| < \pi\ 
\end{aligned} \right.\ \label{eq:9.108}
\end{align}
Die Übertragungsfunktion des Bandpasses kann über die Übertragungsfunktion der Bandsperre ausgedrückt werden.
\begin{align}
G_{BP}(\Omega) &= 1 - G_{BS}(\Omega) \label{eq:9.109}
\end{align}
Beide Amplitudengänge sind in Bild \ref{fig:9.22} dargestellt.

\begin{figure}
  \centering
  \includegraphics[alt={Amplitudengang der idealen Bandsperre und des idealen Bandpasses}]{new media/10_22_amplitudengang_der_idealen_ban.png}
  \caption{Amplitudengang der idealen Bandsperre und des idealen Bandpasses}
  \label{fig:9.22}
\end{figure}

Die Impulsantworten können mit den Rechenregeln der Fourier-Transformation aus den Impulsantworten des idealen Tief- und des idealen Hochpass berechnet werden, sie sind Gegenstand einer Übungsaufgabe.

\subsection{Impulsinvarianter Filterentwurf von FIR-Filtern}\label{sec:impulsinvarianter-filterentwurf-von-fir-filtern}

In Abschnitt \ref{sec:9.3.1} werden die Frequenzgänge und Impulsantworten idealer Filter diskutiert. Sie sind Ausgangspunkt für den impulsinvarianten Entwurf von FIR-Filtern. Der Entwurf wird am Beispiel eines Tiefpass-Filters hergeleitet. Mit den bekannten Frequenzgängen und Impulsantworten von Hochpass, Bandpass und Bandsperre kann das Verfahren auf die anderen Filtertypen übertragen werden.

Der Frequenzgang des zeitkontinuierlichen idealen Tiefpasses mit einer Grenzfrequenz \(ω_{G}\) ist definiert über seinen Amplitudengang
\begin{align}
A_{TP}(\omega) = \left\{ \begin{aligned}
 & 1\text{   f}\text{ü}\text{r }|\omega| < \omega_{G} \\
 & 0\text{  f}\text{ü}\text{r }\omega_{G} < |\omega|
\end{aligned} \right.\ \label{eq:9.110}
\end{align}
und seine lineare Phase
\begin{align}
\varphi_{TP}(\omega) &= e^{- j  \omega  t_{0}} \label{eq:9.111}
\end{align}
Die Impulsantwort des zeitkontinuierlichen idealen Tiefpass-Filters ergibt sich aus den Rechenregeln zur Fourier-Transformation und den angegebenen Korrespondenzen zu
\begin{align}
g(t) &= \frac{\sin\left( \omega_{G}  \left( t - t_{0} \right) \right)}{\pi  \left( t - t_{0} \right)} \label{eq:9.112}
\end{align}
Amplitudengang und Impulsantwort des idealen zeitkontinuierlichen Tiefpass-Filters sind in Bild \ref{fig:9.23} dargestellt.

\begin{figure}
  \centering
  \includegraphics[alt={Impulsantwort des idealen zeitkontinuierlichen Tiefpass-Filters\\
mit einer Grenzfrequenz \(ω_{G} = 1000\) rad/s}]{new media/10_23_impulsantwort_des_idealen_zeit.png}
  \caption{Impulsantwort des idealen zeitkontinuierlichen Tiefpass-Filters\\
mit einer Grenzfrequenz \(ω_{G} = 1000\) rad/s}
  \label{fig:9.23}
\end{figure}

Damit das Filter als zeitdiskretes FIR-Filter mit linearer Phase realisiert werden kann, müssen folgende Bedingungen erfüllt sein:

\begin{itemize}
\item
  Die Impulsantwort \(g[k]\) muss kausal sein.
\item
  Die Impulsantwort \(g[k]\) muss endlich sein.
\item
  Die Impulsantwort \(g[k]\) muss eine geeignete Symmetrie aufweisen.
\end{itemize}

Die ersten beiden Bedingungen werden über das Ausschneiden der Impulsantwort \(g(t)\) über eine Fensterfunktion \(w(t)\) erreicht, die zu Zeitpunkt \(t = 0\) beginnt und zum Zeitpunkt T endet. Wegen der Achsensymmetrie der Impulsantwort kann die Bedingung nach Symmetrie erfüllt werden, wenn die Fensterung des Signals an der Stelle T = 2t\textsubscript{0} endet.
\begin{align}
g_{w}(t) &= g(t)  w(t) = \frac{\sin\left( \omega_{G}  \left( t - t_{0} \right) \right)}{\pi  \left( t - t_{0} \right)}  \left( \sigma(t) - \sigma\left( t - 2  t_{0} \right) \right) \label{eq:9.113}
\end{align}
Damit muss ein zeitdiskretes Filter der Ordnung N bestimmt werden, das an N + 1 Abtastwerten der Impulsantwort dieselben Werte aufweist wie die Impulsantwort des entsprechenden zeitkontinuierlichen Filters. Die Vorgehensweisen zur Bestimmung der Koeffizienten \(b_{m}\) des Filters sind davon abhängig, ob die Filterordnung gerade oder ungerade ist. Beide Varianten werden an einem Beispiel verdeutlicht, bei dem ein Filter mit einer Grenzfrequenz von \(ω_{G} = 1000\) rad/s mit einer Abtastzeit von\\
\(T_{A} = \SI{1}{\milli\second}\) realisiert werden soll.

\subsubsection{Gerade Filterordnung}\label{sec:gerade-filterordnung}

Für eine gerade Filterordnung N werden N + 1 Abtastwerte in das Intervall von 0 ≤ t ≤ 2t\textsubscript{0} gelegt. Damit gilt für die Zeit \(t_{0}\) die Bedingung
\begin{align}
t_{0} &= \frac{N}{2}  T_{A} \label{eq:9.114}
\end{align}
Für das oben beschriebene Beispiel ergibt sich bei einer Filterordnung \(N = 8\) die Zeit \(t_{0}\) = 4T\textsubscript{A} = \SI{4}{\milli\second}, und die Abtastwerte der Impulsantwort berechnen sich für \(k = 0\) \ldots{} 8 zu
\begin{align}
g[ k] &= g\left( k  T_{A} \right) = \frac{\sin\left( \omega_{G}  \left( k  T_{A} - 4  T_{A} \right) \right)}{\pi  \left( k  T_{A} - 4  T_{A} \right)} = \frac{\sin\left( \Omega_{G}  (k - 4) \right)}{\pi  T_{A}  (k - 4)} \label{eq:9.115}
\end{align}
Bild \ref{fig:9.24} stellt die Abtastwerte für das oben beschriebene Beispiel bei einem impulsinvarianten Filterentwurf und einer geraden Filterordnung von \(N = 8\) dar.

\begin{figure}
  \centering
  \includegraphics[alt={Abtastwerte bei einem impulsinvarianten Filterentwurf und gerader Filterordnung \(N = 8\)}]{new media/10_24_abtastwerte_bei_einem_impulsin.png}
  \caption{Abtastwerte bei einem impulsinvarianten Filterentwurf und gerader Filterordnung \(N = 8\)}
  \label{fig:9.24}
\end{figure}

Aus der Impulsantwort \(g[k]\) kann die Übertragungsfunktion \(G(z)\)
\begin{align}
G(z) &= \sum\limits_{k = 0}^{N}{g[ k]  z^{- k}} \label{eq:9.116}
\end{align}
und der Frequenzgang des Filters
\begin{align}
G(\Omega) &= \sum\limits_{k = 0}^{N}{g[ k]  e^{- j  k  \Omega}} \label{eq:9.117}
\end{align}
bestimmt werden.

Bild \ref{fig:9.25} stellt Amplituden- und Phasengang des Filters dar. Zum besseren Vergleich des Ergebnisses mit der Spezifikation werden Amplituden- und Phasengang als Funktion der Kreisfrequenz ω und nicht der normierten Kreisfrequenz Ω dargestellt.

\begin{figure}
  \centering
  \includegraphics[alt={Amplituden- und Phasengang des entworfenen FIR-Filters\\
bei einem impulsinvarianten Filterentwurf und gerader Filterordnung \(N = 8\)}]{new media/10_25_amplituden_und_phasengang_des.png}
  \caption{Amplituden- und Phasengang des entworfenen FIR-Filters\\
bei einem impulsinvarianten Filterentwurf und gerader Filterordnung \(N = 8\)}
  \label{fig:9.25}
\end{figure}

Das Filter erfüllt weitgehend das Spezifikationsziel, an der definierten Grenzfrequenz \(ω_{G} = 1000\) rad/s beträgt der Amplitudengang A(\(ω_{G}\)) = 0.5, das erste Nebenmaximum hat eine Höhe von \(A_{N} = 0.118\). Der Phasengang ist abgesehen von Sprüngen um ∆ϕ = - π an den Stellen, an denen der Amplitudengang Nullstellen aufweist, linear.

Die Differenzengleichung ergibt sich wie bei den Filterentwürfen zuvor aus der Rücktransformation der Übertragungsgleichung \(G(z)\) in den Zeitbereich.

\subsubsection{Ungerade Filterordnung}\label{sec:ungerade-filterordnung}

Auch bei ungerader Filterordnung N werden N + 1 Abtastwerte in das Intervall von 0 ≤ t ≤ 2t\textsubscript{0} gelegt. Damit gilt für die Zeit \(t_{0}\) die Bedingung
\begin{align}
t_{0} &= \frac{N}{2}  T_{A} \label{eq:9.118}
\end{align}
Die Zeit \(t_{0}\) ist damit selbst kein Abtastwert mehr. Für das oben beschriebene Beispiel ergibt sich bei einer Filterordnung \(N = 9\) die Zeit \(t_{0} = 4\).5T\textsubscript{A} = \SI{4.5}{\milli\second}, und die Abtastwerte der Impulsantwort berechnen sich für \(k = 0\) \ldots{} 9 zu
\begin{align}
g[ k] &= g\left( k  T_{A} \right) = \frac{\sin\left( \omega_{G}  \left( k  T_{A} - 4.5  T_{A} \right) \right)}{\pi  \left( k  T_{A} - 4.5  T_{A} \right)} = \frac{\sin\left( \Omega_{G}  (k - 4.5) \right)}{\pi  T_{A}  (k - 4.5)} \label{eq:9.119}
\end{align}
Bild \ref{fig:9.24} stellt die Abtastwerte für das oben beschriebene Beispiel bei einem impulsinvarianten Filterentwurf und einer ungeraden Filterordnung von \(N = 9\) dar.

\begin{figure}
  \centering
  \includegraphics[alt={Abtastwerte bei einem impulsinvarianten Filterentwurf und ungerader Filterordnung \(N = 9\)}]{new media/10_26_abtastwerte_bei_einem_impulsin.png}
  \caption{Abtastwerte bei einem impulsinvarianten Filterentwurf und ungerader Filterordnung \(N = 9\)}
\end{figure}

Aus der Impulsantwort \(g[k]\) kann wieder die Übertragungsfunktion \(G(z)\) und der Frequenzgang des Filters \(G(Ω)\) bestimmt werden.

\begin{figure}
  \centering
  \includegraphics[alt={Vergleich des Amplituden- und Phasengangs der entworfenen FIR-Filter\\
bei einem impulsinvarianten Filterentwurf mit den Ordnungen \(N = 8\) und \(N = 9\)}]{new media/10_27_vergleich_des_amplituden_und_p.png}
  \caption{Vergleich des Amplituden- und Phasengangs der entworfenen FIR-Filter\\
bei einem impulsinvarianten Filterentwurf mit den Ordnungen \(N = 8\) und \(N = 9\)}
\end{figure}

Erwartungsgemäß ist das Filter der Ordnung \(N = 9\) steiler als das Filter der Ordnung \(N = 8\). Die Güte der Approximation des idealen Amplitudengangs wächst mit der Ordnung des Filters. Andererseits führt eine Steigerung der Filterordnung zu einem steileren Phasengang und damit zu einer höheren Gruppenlaufzeit. Die Gruppenlaufzeit \(T_{G}\) steigt mit wachsender Filterordnung linear an.
\begin{align}
T_{G} &= t_{0} = \frac{N}{2}  T_{A} \label{eq:9.120}
\end{align}
Die Differenzengleichung ergibt sich wie bei den Filterentwürfen zuvor aus der Rücktransformation der Übertragungsgleichung \(G(z)\) in den Zeitbereich.

\subsubsection{Zusammenfassung zum impulsinvarianten Entwurf von FIR-Filtern}\label{sec:zusammenfassung-zum-impulsinvarianten-entwurf-von-fir-filtern}

Tabelle \ref{tab:9.6} fasst das Vorgehen zum impulsinvarianten Entwurf von FIR-Filtern zusammen. Dabei wird davon ausgegangen, dass die Filterordnung N und die Abtastzeit \(T_{A}\) vorgegeben sind.

\begin{table}
  \centering
  \caption{Vorgehen zum impulsinvarianten Entwurf von FIR-Filtern}
  \label{tab:9.6}
  \begin{tabular}{ll}
    \toprule
    Schritt & Beschreibung \\
    \midrule
    1 & Bestimmung der Impulsantwort \(g(t)\) des idealen Filters, Beispiel Tiefpass \\
    & \(g(t) = \frac{\sin\left( \omega_{G}  \left( t - t_{0} \right) \right)}{\pi  \left( t - t_{0} \right)}\) \\
    2 & Bestimmung der Zeitverschiebung \(t_{0}\) aus Filterordnung N und Abtastzeit \(T_{A}\) \\
    & \(t_{0} = \frac{N}{2}  T_{A}\) \\
    3 & Bestimmung der N + 1 Werte der Impulsantwort \(g[k]\) \\
    & \(g\lbrack k\rbrack = g\left( k  T_{A} \right)\) \\
    4 & Berechnung der Übertragungsfunktion \\
    & \(G(z) = \sum\limits_{k = 0}^{N}{g\lbrack k\rbrack  z^{- k}}\) \\
    5 & Aufstellen der Differenzengleichung durch Rücktransformation der Gleichung \(G(z)\) in den Zeitbereich \\
    & \(G(z) = \frac{Y(z)}{U(z)}\) \\
    \bottomrule
  \end{tabular}
\end{table}

Alternativ zu dem hier skizzierten Verfahren, bei dem die Gruppenlaufzeit \(T_{G} = t_{0}\) über die Filterordnung N und die Abtastzeit \(T_{A}\) bestimmt wird, ist auch die Vorgabe anderer Randbedingungen wie Gruppenlaufzeit und Filterordnung denkbar. Das Verfahren bleibt in diesen Fällen ähnlich.

\subsection{Approximation des Frequenzgangs}\label{sec:approximation-des-frequenzgangs}

In Abschnitt \ref{sec:9.3.1} werden FIR-Filter ausgehend von der Impulsantwort eines idealen Filters bestimmt. Der Umweg über den Zeitbereich ist nicht zwingend erforderlich. Alternativ können die Koeffizienten des Filters aus der inversen Fourier-Transformation für Signalfolgen bestimmt werden. Nach den Ausführungen in Kapitel \hyperref[spektrum-von-signalfolgen]{7} gilt
\begin{align}
g[ k] &= \frac{1}{2  \pi}  \int_{- \pi}^{\pi}{G(\Omega)  e^{j  k  \Omega}\ d\Omega} \label{eq:9.121}
\end{align}
Dabei ist die Impulsantwort \(g[k]\) nicht zwingend kausal und besitzt außerdem eine unendlich große Länge. Um diese beiden Einschränkungen zu umgehen, wird die Impulsantwort für den Bereich\\
- N/2 ≤ k ≤ N/2 berechnet und anschließend um \(k_{0} = N\)/2 nach rechts verschoben. Da der Wert \(k_{0}\) ein ganzzahliger Wert sein muss, lassen sich mit dem hier beschriebenen Verfahren nur FIR-Filter gerader Ordnung implementieren. Das Verfahren wird an dem Beispiel aus Abschnitt \ref{sec:9.3.1} verdeutlicht, bei dem ein Filter mit einer Grenzfrequenz von \(ω_{G} = 1000\) rad/s mit einer Abtastzeit von \(T_{A} = \SI{1}{\milli\second}\) und eine Ordnung \(N = 8\) realisiert werden soll. Die Grenzfrequenz \(ω_{G}\) entspricht der normierten Grenzfrequenz \(Ω_{G}\) = \(ω_{G}\)\(T_{A} = 1\). Damit berechnen sich die Koeffizienten \(g[k]\) zu
\begin{align}
{g\left[k\right]} &=  {\frac{1}{2 \pi }  \int _{-\pi }^{\pi }G\left(\Omega \right) e^{j k \Omega } {\; }d\Omega  =\frac{1}{2 \pi }  \int _{-\Omega _{G} }^{\Omega _{G} }e^{j k \Omega } {\; }d\Omega  } \\
{} &=  {\frac{1}{2 \pi }  \frac{1}{j k}  \left. e^{j k \Omega } \right|_{-\Omega _{G} }^{\Omega _{G} } =\frac{\sin \left(\Omega _{G}  k\right)}{\pi  k} }
 \label{eq:9.122}
\end{align}
Damit die Impulsantwort kausal wird, wird sie um N/2 = 4 verschoben, und es ergibt sich
\begin{align}
g[ k - 4] &= \frac{\sin\left( \Omega_{G}  (k - 4) \right)}{\pi  (k - 4)} \label{eq:9.123}
\end{align}
Bis auf den Proportionalitätsfaktor 1/\(T_{A}\) entspricht die Impulsantwort der Impulsantwort in Gleichung \eqref{eq:9.115}, die nach dem impulsinvarianten Verfahren bestimmt wird. Beide Verfahren führen damit zu derselben Filterfunktion. Tabelle \ref{tab:9.7} fasst das Vorgehen zum Entwurf von FIR-Filtern über die Approximation des Frequenzgangs zusammen.

\begin{table}
  \centering
  \caption{Vorgehen zum Entwurf von FIR-Filtern über die Approximation des Frequenzgangs}
  \label{tab:9.7}
  \begin{tabular}{ll}
    \toprule
    Schritt & Beschreibung \\
    \midrule
    1 & Definition des idealen Frequenzgangs \(G(Ω)\) \\
    2 & Bestimmung der N + 1 Werte der Impulsantwort \(g[k]\) für - N/2 ≤ k ≤ N/2 \\
    & \(g\lbrack k\rbrack = \frac{1}{2  \pi}  \int_{- \pi}^{\pi}{G(\Omega)  e^{j  k  \Omega}\ d\Omega}\) \\
    3 & Verschiebung der Impulsantwort um \(k_{0} = N\)/2 zur Generierung einer kausalen Impulsantwort \\
    4 & Berechnung der Übertragungsfunktion \\
    & \(G(z) = \sum\limits_{k = 0}^{N}{g\lbrack k\rbrack  z^{- k}}\) \\
    5 & Aufstellen der Differenzengleichung durch Rücktransformation der Gleichung \(G(z)\) in den Zeitbereich \\
    & \(G(z) = \frac{Y(z)}{U(z)}\) \\
    \bottomrule
  \end{tabular}
\end{table}

\subsection{Fensterfunktionen zur Verbesserung der Filtercharakteristik}\label{sec:fensterfunktionen-zur-verbesserung-der-filtercharakteristik}

Bei der Berechnung der Impulsantwort \(g_{w}\){[}k{]} wird die Impulsantwort \(g(t)\) abgetastet und mit einer Fensterfunktion \(w[k]\) multipliziert.
\begin{align}
g_{w}[ k] &= g[ k]  w[ k] \label{eq:9.124}
\end{align}
Nach den Rechenregeln zur Fourier-Transformation für Signalfolgen in Abschnitt \ref{sec:7.2.9} ergibt sich bei Multiplikation zweier Signalfolgen im Zeitbereich eine Faltung der jeweiligen Spektren im Frequenzbereich.
\begin{align}
G_{w}(\Omega) &= G(\Omega)*W(\Omega) = \frac{1}{2  \pi}  \int_{- \pi}^{\pi}{G(\Theta)  W(\Omega - \Theta)\ d\Theta} \label{eq:9.125}
\end{align}
Durch die Faltung des Spektrums \(G(Ω)\) mit dem Spektrum der Fensterfunktion \(W(Ω)\) wird der Frequenzgang verändert. Wäre das Fenster \(w[k]\) konstant und unendlich lang, ergäbe sich für das Spektrum \(W(Ω)\) eine periodische Impulsfolge δ(Ω - k2π). Im Frequenzbereich - π ≤ Ω ≤ π liegt der Impuls an der Stelle \(Ω = 0\). Die Faltung eines Spektrums mit einem Impuls an der Stelle \(Ω = 0\) verändert das Spektrum nicht, sodass ein Filter mit einer unendlich langen Impulsantwort \(g_{w}\){[}k{]} zu dem eigentlich gewünschten Frequenzgang \(G(Ω)\) führen würde. Mit dieser Forderung ist allerdings bei der Implementierung ein unendlich hoher Rechenaufwand verbunden, sodass dieser Ansatz ausscheidet. Deshalb ist das Ziel, eine Fensterung zu finden, die bei einer endlichen Länge der Fensterfunktion \(w[k]\) ein möglichst schmales Spektrum besitzt, das der Impulsfunktion möglichst nahekommt.

Zur Diskussion unterschiedlicher Fensterfunktionen ist in Bild \ref{fig:9.28} das Spektrum des Rechteck-Fensters dargestellt.

\begin{figure}
  \centering
  \includegraphics[alt={Kriterien zur Bewertung von Spektrum von Fensterfunktionen}]{new media/10_28_kriterien_zur_bewertung_von_sp.png}
  \caption{Kriterien zur Bewertung von Spektrum von Fensterfunktionen}
  \label{fig:9.28}
\end{figure}

Das Spektrum besteht aus einem Hauptschwinger (Main-Lobe) und den Nebenschwingern (Side-Lobes). Die Spektren realer Fensterfunktionen haben eine endliche Breite des Hauptmaximums, und es existieren Nebenmaxima mit endlich hohen Amplituden. Zur Charakterisierung von Fensterfunktionen werden folgende Bewertungskriterien verwendet.

\begin{itemize}
\item
  Breite des Hauptmaximums\\
  Die Breite des Hauptmaximums ist der Frequenzbereich zwischen \(Ω = 0\) und dem ersten Nulldurchgang des Spektrums \(Ω_{S}\). Eine gute Fensterfunktion hat eine geringe Breite des Hauptmaximums, da dadurch bei der Faltung das Spektrum der gefensterten Funktion weniger stark verbreitert wird.
\item
  Relative Amplitude des ersten Nebenmaximums\\
  Die relative Amplitude des Nebenmaximums ist das Verhältnis von Maximum des Nebenmaximums \(A_{SL}\) zu dem Hauptmaximum \(A_{ML}\). Die relative Amplitude wird in dB angegeben.
\end{itemize}
\begin{align}
a_{REL} &= 20  \log\left( \frac{A_{SL}}{A_{ML}} \right) \label{eq:9.126}
\end{align}
Zielwert ist eine geringe Höhe des Nebenmaximums, also eine kleine relative Amplitude des Nebenmaximums.

Mit der Bestimmung des idealen Fensters beschäftigen sich viele Veröffentlichungen. Es zeigt sich, dass ein ideales Fenster nicht existiert, sondern dass ein Kompromiss zwischen endlicher Breite und relativer Amplitude des Nebenmaximums eingegangen werden muss. Einige Fenster werden im Folgenden beschrieben. Dabei wird zur Vereinfachung der Darstellung von einer Fensterfunktion ausgegangen, die symmetrisch zum Zeitpunkt \(k = 0\) liegt. Diese Fenster sind nicht kausal, können aber durch eine Zeitverschiebung in kausale Fenster überführt werden.

\subsubsection{Rechteckfenster}\label{sec:rechteckfenster}

Das Rechteckfenster wird bei gerader Filterordnung N im Zeitbereich über die Gleichung
\begin{align}
w_{REC}[ k] &= \sigma\left( k + \frac{N}{2} \right) - \sigma\left( k - \frac{N}{2} - 1 \right) \label{eq:9.127}
\end{align}
beschrieben und hat als Spektrum den sogenannten Dirichlet-Kern
\begin{align}
W_{REC}(\Omega) &= \frac{\sin\left( \frac{N  \Omega}{2} \right)}{\sin\left( \frac{\Omega}{2} \right)} \label{eq:9.128}
\end{align}
Das Rechteckfenster und der normierte Amplitudengang sind in Bild \ref{fig:9.29} dargestellt.

\begin{figure}
  \centering
  \includegraphics[alt={Rechteckfenster und Spektrum des Rechteckfensters für \(N = 20\)}]{new media/10_29_rechteckfenster_und_spektrum_d.png}
  \caption{Rechteckfenster und Spektrum des Rechteckfensters für \(N = 20\)}
  \label{fig:9.29}
\end{figure}

Das Hauptmaximum des Rechteckfensters hat eine Breite von \(Ω_{S} = 2\)·π/N. Die relative Amplitude des Nebenmaximums beträgt \(a_{REL}\) = - \SI{13}{\decibel}.

\subsubsection{Dreieck- oder Bartlett-Fenster}\label{sec:dreieck--oder-bartlett-fenster}

Das Dreieckfenster wird im Zeitbereich über die Gleichung
\begin{align}
w_{TRI}[ k] &= \frac{2}{N}  \left( \left( k + \frac{N}{2} \right)  \sigma\left( k + \frac{N}{2} \right) - 2  k  \sigma[ k] + \left( k - \frac{N}{2} \right)  \sigma\left( k - \frac{N}{2} \right) \right) \label{eq:9.129}
\end{align}
beschrieben und hat das Spektrum
\begin{align}
W_{TRI}(\Omega) &= \frac{\sin^{2}\left( \frac{N  \Omega}{4} \right)}{\sin^{2}\left( \frac{\Omega}{2} \right)} \label{eq:9.130}
\end{align}
Das Dreieckfenster und der normierte Amplitudengang sind in Bild \ref{fig:9.30} dargestellt.

\begin{figure}
  \centering
  \includegraphics[alt={Dreieckfenster und Spektrum des Dreieckfensters für \(N = 20\)}]{new media/10_30_dreieckfenster_und_spektrum_de.png}
  \caption{Dreieckfenster und Spektrum des Dreieckfensters für \(N = 20\)}
  \label{fig:9.30}
\end{figure}

Das Hauptmaximum des Dreieckfensters hat eine Breite von \(Ω_{S} = 4\)·π/N. Die relative Amplitude des Nebenmaximums beträgt \(a_{REL}\) = - \SI{25}{\decibel}.

\subsubsection{Hann-Fenster}\label{sec:hann-fenster}

Auch das Hann-Fenster weist am Beginn und am Ende der Beobachtungszeit keine Sprünge auf, sondern geht an diesen Punkten gegen den Wert 0. Der stetige Übergang wird mit einer Kosinusfunktion beschrieben.
\begin{align}
w_{HAN}[ k] &= \left( 0.5 + 0.5  \cos\left( \frac{2  \pi  k}{N} \right) \right)  \left( \sigma\left( k + \frac{N}{2} \right) - \sigma\left( k - \frac{N}{2} - 1 \right) \right) \label{eq:9.131}
\end{align}
Die Fensterfunktion ergibt sich aus dem Produkt von Rechteckfenster und der Kosinusfunktion. Das Spektrum des Hann-Fensters kann über die Faltungsoperation berechnet werden. Es ergibt sich zu
\begin{align}
W_{HAN}(\Omega) &= \frac{1}{4}  \frac{\sin\left( \frac{N  \left( \Omega + \frac{2  \pi}{N} \right)}{2} \right)}{\sin\left( \frac{\Omega + \frac{2  \pi}{N}}{2} \right)} + \frac{1}{2}  \frac{\sin\left( \frac{N  \Omega}{2} \right)}{\sin\left( \frac{\Omega}{2} \right)} + \frac{1}{4}  \frac{\sin\left( \frac{N  \left( \Omega - \frac{2  \pi}{N} \right)}{2} \right)}{\sin\left( \frac{\Omega - \frac{2  \pi}{N}}{2} \right)} \label{eq:9.132}
\end{align}
Das Hann-Fenster und sein Spektrum sind in Bild \ref{fig:9.31} dargestellt.

\begin{figure}
  \centering
  \includegraphics[alt={Hann-Fenster und Spektrum des Hann-Fensters für \(N = 20\)}]{new media/10_31_hannfenster_und_spektrum_des_h.png}
  \caption{Hann-Fenster und Spektrum des Hann-Fensters für \(N = 20\)}
  \label{fig:9.31}
\end{figure}

Das Hauptmaximum des Hann-Fensters hat eine Breite von \(Ω_{S} = 4\)·π/N und ist damit doppelt so breit wie bei dem Rechteckfenster. Die relative Amplitude des Nebenmaximums ist mit \(a_{REL}\) = - \SI{32}{\decibel} deutlich größer als bei dem Rechteckfenster.

\subsubsection{Hamming-Fenster}\label{sec:hamming-fenster}

Hamming- und Hann-Fenster unterscheiden sich nur geringfügig in der Definition.
\begin{align}
w_{HAM}[ k] &= \left( 0.54 + 0.46  \cos\left( \frac{2  \pi  k}{N} \right) \right)  \left( \sigma\left( k + \frac{N}{2} \right) - \sigma\left( k - \frac{N}{2} - 1 \right) \right) \label{eq:9.133}
\end{align}
Auch die Berechnung des Spektrums der beiden Fensterfunktionen ist ähnlich. Für das Hamming-Fenster ergibt sich das Spektrum
\begin{align}
W_{HAM}(\Omega) &= 0.23  \frac{\sin\left( \frac{N  \left( \Omega + \frac{2  \pi}{N} \right)}{2} \right)}{\sin\left( \frac{\Omega + \frac{2  \pi}{N}}{2} \right)} + 0.54  \frac{\sin\left( \frac{N  \Omega}{2} \right)}{\sin\left( \frac{\Omega}{2} \right)} + 0.23  \frac{\sin\left( \frac{N  \left( \Omega - \frac{2  \pi}{N} \right)}{2} \right)}{\sin\left( \frac{\Omega - \frac{2  \pi}{N}}{2} \right)} \label{eq:9.134}
\end{align}
Das Hamming-Fenster und sein Spektrum sind in Bild \ref{fig:9.32} dargestellt.

\begin{figure}
  \centering
  \includegraphics[alt={Hamming-Fenster und Spektrum des Hamming-Fensters für \(N = 20\)}]{new media/10_32_hammingfenster_und_spektrum_de.png}
  \caption{Hamming-Fenster und Spektrum des Hamming-Fensters für \(N = 20\)}
  \label{fig:9.32}
\end{figure}

Das Hauptmaximum des Hamming-Fensters hat eine Breite von \(Ω_{S} = 4\)·π/N. Die relative Amplitude des Nebenmaximums beträgt mit \(a_{REL}\) = - \SI{42}{\decibel} und ist nochmals deutlich größer als bei dem Hann-Fenster.

\subsubsection{Vergleich der Fensterfunktionen}\label{sec:vergleich-der-fensterfunktionen}

Tabelle \ref{tab:9.8} stellt die berechneten Kennwerte der unterschiedlichen Fenster zusammen.

\begin{table}
  \centering
  \caption{Tabellarischer Vergleich der Kennwerte unterschiedlicher Fensterfunktionen}
  \label{tab:9.8}
  \begin{tabular}{lll}
    \toprule
    Fenster- 
Funktion
 & 
Breite des 
Hauptmaximums \(Ω_{S}\)
 & 
Relative Amplitude 
der Nebenmaxima \(a_{REL}\) \\
    \midrule
    Rechteck & 2·π / N & - \SI{13}{\decibel} \\
    Dreieck & 4·π / N & - \SI{25}{\decibel} \\
    Hann & 4·π / N & - \SI{32}{\decibel} \\
    Hamming & 4·π / N & - \SI{42}{\decibel} \\
    \bottomrule
  \end{tabular}
\end{table}

Die Daten zeigen, dass es eine ideale Fensterfunktion nicht gibt. Das Spektrum des Rechteckfensters ist schmal, aber es weist vergleichsweise hohe Nebenmaxima auf. Demgegenüber haben die Spektren von Hann- und Hamming-Fenster relativ kleine Nebenmaxima, die Hauptmaxima sind aber breiter als bei dem Rechteckfenster. Für praktische Anwendungen bietet das Hamming-Fenster einen guten Kompromiss zwischen Breite des Hauptmaximums und relativer Amplitude des Nebenmaximums.

\begin{example}{Entwicklung eines Tiefpass-Filters mit Hamming-Fenster}
\label{beispiel-entwicklung-eines-tiefpass-filters-mit-hamming-fenster}

Das Verfahren wird an dem Beispiel eines Tiefpass-Filters mit einer Grenzfrequenz von \(ω_{G} = 1000\) rad/s und einer Abtastzeit von \(T_{A} = \SI{1}{\milli\second}\) angewendet. Die Ordnung des Filters wird auf \(N = 12\) festgelegt. Damit ergibt sich eine Gruppenlaufzeit \(T_{G} = t_{0}\) = \SI{6}{\milli\second}.
\begin{align}
g[ k] &= \frac{\sin\left( \omega_{G}  \left( k  T_{A} - 6  T_{A} \right) \right)}{\pi  \left( k  T_{A} - 6  T_{A} \right)} = \frac{\sin\left( \Omega_{G}  (k - 6) \right)}{\pi  T_{A}  (k - 6)} \label{eq:9.135}
\end{align}
Die Impulsantwort \(g[k]\) wird mit dem um \(k = 6\) nach rechts verschobenen Hamming-Fenster multipliziert, und es ergibt sich die gefensterte Impulsantwort
\begin{align}
g_{w}[ k] &= \frac{\sin\left( \Omega_{G}  (k - 6) \right)}{\pi  T_{A}  (k - 6)}  \left( 0.54 + 0.46  \cos\left( \frac{2  \pi  (k - 6)}{N} \right) \right) \label{eq:9.136}
\end{align}
Bild \ref{fig:9.33} vergleicht die Impulsantworten für die Filterentwürfe mit Rechteck- und Hamming-Fenster miteinander.

\begin{fixedfigure}
  \includegraphics[alt={Impulsantworten der Filterentwürfe mit Rechteck- und Hamming-Fenster im Vergleich,\\
Tiefpassfilter mit einer Grenzfrequenz \(ω_{G} = 1\) krad/s, \(N = 12\)}]{new media/10_33_impulsantworten_der_filterentw.png}
  \captionof{figure}{Impulsantworten der Filterentwürfe mit Rechteck- und Hamming-Fenster im Vergleich,\\
Tiefpassfilter mit einer Grenzfrequenz \(ω_{G} = 1\) krad/s, \(N = 12\)}
  \label{fig:9.33}
\end{fixedfigure}

Die überlagerte Fensterfunktion des Hamming-Fensters verändert die Filterkoeffizienten. Anfangswert \(g_{w}\){[}0{]} und Endwert \(g_{w}\){[}12{]} werden dadurch zu null. Aus der Impulsantwort kann direkt die Übertragungsfunktion bestimmt werden. Das mit dem Hamming-Fenster entwickelte System hat die Übertragungsfunktion
\begin{align}
G_{W}(z) &= \sum\limits_{k = 0}^{12}{g_{w}[ k]  z^{- k}} = \sum\limits_{k = 0}^{12}{\frac{\sin\left( \Omega_{G}  (k - 6) \right)}{\pi  T_{A}  (k - 6)}  \left( 0.54 + 0.46  \cos\left( \frac{2  \pi  (k - 6)}{N} \right) \right)  z^{- k}} \label{eq:9.137}
\end{align}
Da FIR-Systeme immer stabil sind, ergibt sich der Frequenzgang zu
\begin{align}
G_{W}(\Omega) &= \sum\limits_{k = 0}^{12}{g_{w}[ k]  e^{- j  \Omega  k}} = \sum\limits_{k = 0}^{12}{\frac{\sin\left( \Omega_{G}  (k - 6) \right)}{\pi  T_{A}  (k - 6)}  \left( 0.54 + 0.46  \cos\left( \frac{2  \pi  (k - 6)}{N} \right) \right)  e^{- j  \Omega  k}} \label{eq:9.138}
\end{align}
Die zu den Impulsantworten gehörenden Frequenzgänge werden in Bild \ref{fig:9.34} miteinander verglichen.

\begin{fixedfigure}
  \includegraphics[alt={Filterentwurf mit Rechteck- und Hamming-Fenster im Vergleich,\\
Tiefpassfilter mit einer Grenzfrequenz \(ω_{G} = 1\) krad/s, \(N = 12\)}]{new media/10_34_filterentwurf_mit_rechteck_und.png}
  \captionof{figure}{Filterentwurf mit Rechteck- und Hamming-Fenster im Vergleich,\\
Tiefpassfilter mit einer Grenzfrequenz \(ω_{G} = 1\) krad/s, \(N = 12\)}
  \label{fig:9.34}
\end{fixedfigure}

Es wird deutlich, dass durch den Einsatz des Hamming-Fensters die Welligkeit deutlich kleiner ist als bei dem Rechteckfenster. Dies betrifft sowohl den Durchlass als auch den Sperrbereich des Filters. Dafür weist das mit Rechteckfenster entworfene Filter einen deutlich steileren Übergang zwischen Durchlass- und Sperrbereich auf. Auch die Phasengänge unterscheiden sich, sind aber beide linear. Die Gruppenlaufzeit ist mit \(T_{G} = \SI{6}{\milli\second}\) für beide Entwürfe identisch.

\end{example}

\section{Numerischer Entwurf von FIR-Filtern}\label{sec:9.4}

Bei dem analytischen Entwurf von FIR-Filtern wird typischerweise die in Abschnitt \ref{sec:9.3} beschriebene Methode der Fensterung verwendet. Oftmals wird jedoch ein Filter benötigt, das hinsichtlich definierter Merkmale optimal ist. Bei dem hier beschriebenen Filterentwurf sind das die Merkmale

\begin{itemize}
\item
  Gleichwelligkeit im Durchlassbereich und im Sperrbereich
\item
  Minimale Filter-Ordnung
\item
  Minimierung des maximalen Fehlers
\end{itemize}

Insbesondere die Gleichwelligkeit ist eine Eigenschaft, die bei den analytischen Entwurfsverfahren für FIR-Filter nicht berücksichtigt wird. Bild \ref{fig:9.35} zeigt das Ergebnis eines Tiefpassentwurfs mit dem sogenannten Remez-Entwurf. Im Durchlass- und Sperrbereich verläuft der Frequenzgang gleichwellig in den Toleranzbändern. Die Toleranzbänder in Durchlass- und Sperrbereich können unabhängig voneinander gewählt werden.

\begin{figure}
  \centering
  \includegraphics[alt={Ergebnis für einen Filter nach dem Remez-Entwurf}]{new media/10_35_ergebnis_fr_einen_filter_nach.png}
  \caption{Ergebnis für einen Filter nach dem Remez-Entwurf}
  \label{fig:9.35}
\end{figure}

Das in diesem Abschnitt beschriebene numerische Verfahren zum Entwurf von FIR-Filtern ist in der Literatur in unterschiedlichen Varianten und Namen zu finden, wie Equiripple-Design oder Remez-Entwurf. Zur Approximation des Frequenzgangs werden dabei sogenannte Tschebyscheff-Polynome verwendet, die auf harmonischen Funktionen basieren. Parks und McClellan haben das in diesem Kapitel beschriebene Verfahren durch den Einsatz der Lagrange-Funktion zur Approximation weiterentwickelt. Der Übersicht halber konzentriert sich diese Beschreibung auf den Remez-Entwurf eines Tiefpasses, der auf dem Alternantentheorem und dem Remez-Austauschalgorithmus basiert.

\subsection{Alternanten und das Alternantentheorem}\label{sec:9.4.1}

Eine Alternante ist ein Berührungspunkt eines Polynoms mit einem Toleranzbereich. An den Berührungspunkten muss das Polynom einen Extremwert aufweisen oder auf einem Eckpunkt des Toleranzbereichs liegen. Außerdem müssen die Punkte um eine Grundlinie alternieren.

\begin{figure}
  \centering
  \includegraphics[alt={Toleranzbereich und Polynom mit Alternanten}]{new media/10_36_toleranzbereich_und_polynom_mi.png}
  \caption{Toleranzbereich und Polynom mit Alternanten}
  \label{fig:9.36}
\end{figure}

Bild \ref{fig:9.36} zeigt ein Polynom und einen Toleranzbereich mit Grundlinie. Die markierten Punkte stellen Alternanten dar. Das eingezeichnete Polynom besitzt an diesen Punkten eine Ableitung von null und damit einen Extremwert, oder es verläuft durch einen Eckpunkt des Toleranzbereichs. Außerdem alternieren die Punkte um die eingezeichnete Grundlinie.

Der Remez-Entwurf nutzt das Alternantentheorem für einen Filterentwurfs. Der Amplitudengang des Filters wird beim Remez-Entwurf über Tschebyscheff-Polynome der Ordnung N beschrieben.
\begin{align}
A(\Omega) &= \sum\limits_{n = 0}^{N}{a_{n}  \cos^{n}(\Omega)} \label{eq:9.139}
\end{align}
Die Abweichung des Tschebyscheff-Polynoms \(A(Ω)\) von einer definierten Zielfunktion \(H(Ω)\) wird mit einem Gewichtungsfaktor \(W(Ω)\) gewichtet, der in unterschiedlichen Teilbereichen unterschiedlich groß sein kann. Es ergibt sich die Fehlerfunktion
\begin{align}
E(\Omega) &= W(\Omega)  \left( H(\Omega) - A(\Omega) \right) \label{eq:9.140}
\end{align}
Nach dem Alternantentheorem ist der maximale absolute Fehler minimal, wenn die Fehlerfunktion \(E(Ω)\) mindestens N + 2 Alternanten aufweist.

Die vorgegebene Zielfunktion \(H(Ω)\) besteht aus konstanten Bereichen. Typischerweise weist sie im Durchlassbereich den Wert 1 und im Sperrbereich den Wert null auf, im Übergangsbereich zwischen Durchlass- und Sperrbereich ist sie nicht definiert. In Bild \ref{fig:9.37} links ist die Zielfunktion eingezeichnet. Die abschnittsweise definierten Geraden werden von ihrem jeweiligen Toleranzband mit einem Fehler von ± \(δ_{D}\) beziehungsweise ± \(δ_{S}\) umschlossen. Die Abschnitte und Fehler definieren den Bereich, auf den das Alternantentheorem angewendet wird. Nach dem Alternantentheorem müssen mindestens N + 2 lokale Minima und Maxima an den Stellen
\begin{align}
\Omega_{1} < \Omega_{2} < ... < \Omega_{N + 2} \label{eq:9.141}
\end{align}
existieren, bei denen die Fehlerfunktion ein wechselndes Vorzeichen hat.
\begin{align}
\ E\left( \Omega_{n} \right) &= - E\left( \Omega_{n + 1} \right) \label{eq:9.142}
\end{align}
Ist diese Bedingung erfüllt, ist der maximale absolute Fehler minimiert. Bild \ref{fig:9.37} zeigt den Amplitudengangs eines Tiefpasses, der mit einem Polynom \(A(Ω)\) der Ordnung \(N = 4\) realisiert wurde, und die dazugehörige Fehlerfunktion \(E(Ω)\). An den Toleranzgrenzen des Tiefpasses ist zu erkennen, dass das Polynom N + 2 = 6 Alternanten aufweist. Das Polynom \(A(Ω)\) liegt gleichwellig in den jeweiligen Toleranzbereichen. Nach dem Alternantentheorem ist der kleinstmögliche Fehler gefunden.

\begin{figure}
  \centering
  \includegraphics[alt={Beschreibung eines Amplitudengangs mit einem Polynom der Ordnung \(N = 4\) und Fehlerfunktion}]{new media/10_37_beschreibung_eines_amplitudeng.png}
  \caption{Beschreibung eines Amplitudengangs mit einem Polynom der Ordnung \(N = 4\) und Fehlerfunktion}
  \label{fig:9.37}
\end{figure}

\subsection{Remez-Austauschalgorithmus}\label{sec:remez-austauschalgorithmus}

Das in Abschnitt \ref{sec:9.4.1} beschriebene Alternantentheorem liefert eine Bedingung für einen minimalen Fehler. Auch wenn das Theorem nicht explizit angibt, wie das optimale Polynom zu finden ist, dient es als Basis für einen Algorithmus zur Ermittlung eines solchen Polynoms. Aus dem Alternantentheorem geht hervor, dass das optimale Polynom N-ter Ordnung eine Fehlerfunktion \(E(Ω)\) besitzt.
\begin{align}
E(\Omega) &= W(\Omega)  \left( H(\Omega) - A(\Omega) \right) \label{eq:9.143}
\end{align}
Mit Forderung nach gleichwelligem Verhalten des Polynoms in den vorgegebenen Teilbereichen kann ein Gleichungssystem aufgestellt werden. An den Stellen der gesuchten Alternanten \(Ω_{n}\) weist das Polynom den größten begangenen Fehler δ auf. Dieser alterniert um die vorgegebene Zielfunktion \(H(Ω)\). Für die Alternanten \(Ω_{n}\) mit den Indizes \(n = 1\) \ldots{} N + 2 kann jeweils eine Gleichung
\begin{align}
\delta  ( - 1)^{n + 1} &= W\left( \Omega_{n} \right)  \left( H\left( \Omega_{n} \right) - A\left( \Omega_{n} \right) \right) \label{eq:9.144}
\end{align}
aufgestellt werden. Durch die Gewichtung \(W(Ω)\) kann der Fehler \(\delta\) in allen Teilbereichen als gleich groß angenommen werden. Mit der Definition der Tschebyscheff-Polynome N-ter Ordnung
\begin{align}
A(\Omega) &= \sum\limits_{k = 0}^{N}{a_{k}  \cos^{k}(\Omega)} \label{eq:9.145}
\end{align}
ergibt sich für jede Alternante \(Ω_{n}\) eine Gleichung der Form
\begin{align}
\delta  ( - 1)^{n + 1} &= W\left( \Omega_{n} \right)  \left( H\left( \Omega_{n} \right) - A\left( \Omega_{n} \right) \right) = W\left( \Omega_{n} \right)  \left( H\left( \Omega_{n} \right) - \sum\limits_{k = 0}^{N}{a_{k}  \cos^{k}\left( \Omega_{n} \right)} \right) \label{eq:9.146}
\end{align}
beziehungsweise
\begin{align}
\sum\limits_{k = 0}^{N}{a_{k}  \cos^{k}\left( \Omega_{n} \right)} + \frac{\delta  ( - 1)^{n + 1}}{W\left( \Omega_{n} \right)} &= H\left( \Omega_{n} \right) \label{eq:9.146}
\end{align}
Unter Berücksichtigung der Stützstellen
\begin{align}
\Omega_{1} < \Omega_{2} < ... < \Omega_{N + 2} \label{eq:9.147}
\end{align}
kann ein Gleichungssystem in Matrizenschreibweise erstellt werden.
\begin{align}
\begin{pmatrix}
1 & \cos\left( \Omega_{1} \right) & \cos^{2}\left( \Omega_{1} \right) & s & \cos^{N}\left( \Omega_{1} \right) & \frac{1}{W\left( \Omega_{1} \right)} \\
1 & \cos\left( \Omega_{2} \right) & \cos^{2}\left( \Omega_{2} \right) & s & \cos^{N}\left( \Omega_{2} \right) & \frac{- 1}{W\left( \Omega_{2} \right)} \\
 \vdots & \vdots & \vdots & \ddots & \vdots & \vdots \\
1 & \cos\left( \Omega_{N + 2} \right) & \cos^{2}\left( \Omega_{N + 2} \right) & s & \cos^{N}\left( \Omega_{N + 2} \right) & \frac{( - 1)^{N + 2 + 1}}{W\left( \Omega_{N + 2} \right)}
\end{pmatrix}  \begin{pmatrix}
a_{0} \\
\begin{aligned}
 & a_{1} \\
 & a_{2}
\end{aligned} \\
\begin{aligned}
 & \vdots \\
 & a_{N}
\end{aligned} \\
\delta
\end{pmatrix} = \begin{pmatrix}
H\left( \Omega_{1} \right) \\
H\left( \Omega_{2} \right) \\
 \vdots \\
H\left( \Omega_{N + 2} \right)
\end{pmatrix} \label{eq:9.148}
\end{align}
Das Gleichungssystem kann nach den gesuchten Koeffizient \(a_{k}\) und der maximale Feher δ aufgelöst werden. Da weder die Frequenzstützstellen \(Ω_{n}\) noch die Polynomkoeffizienten \(a_{k}\) bekannt sind, erfolgt die Berechnung der Koeffizienten iterativ.

\subsubsection{Schätzen der initialen Positionen für die Alternanten}\label{sec:schuxe4tzen-der-initialen-positionen-fuxfcr-die-alternanten}

Das Verfahren beginnt mit dem Schätzen von N + 2 Frequenzstützstellen \(Ω_{n}\) in den Bereichen, in denen das Alternantentheorem angewendet werden soll. In Bild \ref{fig:9.38} ist der erste Iterationsschritt für ein Tschebyscheff-Polynom sechsten Grades dargestellt. Für die initiale Berechnung werden acht Frequenzstützstellen \(Ω_{n}\) gleichmäßig auf die Bereiche verteilt. Es ergeben sich vier Alternanten, die auf den Rändern der Toleranzbereiche liegen. Das Alternantentheorem ist nicht erfüllt, das Minimum des Fehlers ist noch nicht erreicht.

\begin{figure}
  \centering
  \includegraphics[alt={Initialisierung der Stützstellen \(Ω_{n}\) beim Remez-Algorithmus und Bestimmung der Extremwerte}]{new media/10_38_initialisierung_der_sttzstelle.png}
  \caption{Initialisierung der Stützstellen \(Ω_{n}\) beim Remez-Algorithmus und Bestimmung der Extremwerte}
  \label{fig:9.38}
\end{figure}

\subsubsection{\texorpdfstring{Austausch der Stützstellen \(Ω_{n}\)}{Austausch der Stützstellen Ωn}}\label{austausch-der-stuxfctzstellen-ux3c9n}

Für den zweiten Iterationsschritt werden die zuvor geschätzten Frequenzen ausgetauscht. Die neuen Werte ergeben sich aus den Extremstellen des Polynoms sowie den beiden Eckfrequenzen \(Ω_{4} = Ω_{D}\) und \(Ω_{5} = Ω_{S}\). Nach Austauschen der Frequenzen wird das Gleichungssystem aus \eqref{eq:9.148} wieder aufgestellt, ausgerechnet und die Alternanten neu bestimmt. Dieses Schema wird so lange wiederholt, bis das Polynom acht Alternanten aufweist und das Alternantentheorem erfüllt ist. Das Ergebnis ist in Bild \ref{fig:9.39} dargestellt.

\begin{figure}
  \centering
  \includegraphics[alt={Polynom nach der Erfüllung des Alternantentheorems}]{new media/10_39_polynom_nach_der_erfllung_des.png}
  \caption{Polynom nach der Erfüllung des Alternantentheorems}
  \label{fig:9.39}
\end{figure}

Ist das Alternantentheorem erfüllt, ist der maximale Fehler δ bekannt und das Tschebyscheff-Polynom ist definiert.

\subsubsection{Ablauf des Remez-Austauschalgorithmus als Flussdiagramm}\label{sec:ablauf-des-remez-austauschalgorithmus-als-flussdiagramm}

Das Flussdiagramm in Bild \ref{fig:9.40} fasst die Vorgehensweise für den Remez-Austauschalgorithmus zusammen.

\begin{figure}
  \centering
  \includesvg{new media/10_40_flussdiagramm_zum_remezaustaus.svg}
  \caption{Flussdiagramm zum Remez-Austauschalgorithmus}
  \label{fig:9.40}
\end{figure}

\subsection{Remez-Filterentwurf}\label{sec:remez-filterentwurf}

Der Remez-Entwurf ist ein Filterentwurf, der auf dem Remez-Austauschalgorithmus basiert. Ausgangspunkt ist ein Toleranzschema, das von dem Filter erfüllt werden muss.

\begin{figure}
  \centering
  \includegraphics[alt={Toleranzschema am Beispiel eines Tiefpasses}]{new media/10_41_toleranzschema_am_beispiel_ein.png}
  \caption{Toleranzschema am Beispiel eines Tiefpasses}
\end{figure}

Der Durchlassbereich liegt beim Tiefpass in Bereich 0 ≤ Ω ≤ \(Ω_{D}\). In dem Bereich weist die Zielfunktion den Wert 1 auf. Die Zielfunktion wird von einem Toleranzbereich von ± \(δ_{D}\) umschlossen. Der Sperrbereich umfasst den Bereich \(Ω_{S}\) ≤ Ω ≤ π und weist ein Toleranzband von ± \(δ_{S}\) auf. Das Toleranzschema spezifiziert den Verlauf des Amplitudengangs nur in den Teilbereichen des Durchlass- und Sperrbereichs. In dem Bereich zwischen den beiden Eckfrequenzen \(Ω_{D}\) und \(Ω_{S}\) kann die Approximationsfunktion eine beliebige Form annehmen. Es wird sich aber zeigen, dass das approximierte Polynom wegen der Kosinusfunktionen als Basis ein gutmütiges Verhalten aufweisen wird.

Die beiden Fehlergrenzen \(δ_{D}\) und \(δ_{S}\) ermöglichen unterschiedliche Fehlertoleranzen in Durchlass- und Sperrbereich. Sie führen zu einer Gewichtung \(W(Ω)\) des Fehlers.
\begin{align}
W(\Omega) = \left\{ \begin{aligned}
 & \frac{\delta_{S}}{\delta_{D}}für0 \leq \Omega \leq \Omega_{D} \\
 & 1für\Omega_{S} \leq \Omega \leq \pi
\end{aligned} \right.\ \label{eq:9.149}
\end{align}
\subsubsection{Abschätzung der initialen Filterordnung}\label{sec:abschuxe4tzung-der-initialen-filterordnung}

Kaiser hat eine Gleichung für die Abschätzung der erforderlichen Ordnung N des Tschebyscheff-Polynoms hergeleitet (Oppenheim, Zeitdiskrete Signalverarbeitung 2004). Für Tiefpässe führt die Abschätzung zu einer Filterordnung von
\begin{align}
N \approx \frac{- 10  \log_{10}\left( \delta_{D}  \delta_{S} \right) - 13}{2  2.324  \left( \Omega_{S} - \Omega_{D} \right)} \label{eq:9.150}
\end{align}
\subsubsection{Anwendung des Remez-Austauschalgorithmus}\label{sec:anwendung-des-remez-austauschalgorithmus}

Mit diesen Vorarbeiten liegen alle Informationen zur Anwendung des Remez-Austauschalgorithmus vor. Nach seiner Anwendung ist bekannt, welcher Fehler δ mit der vorgegebenen Ordnung N erreicht werden kann. Eine Prüfung der Toleranzvorgaben
\begin{align}
\delta_{S} \geq \delta \label{eq:9.151}
\end{align}
und
\begin{align}
\delta_{D} \geq \frac{\delta  \delta_{D}}{\delta_{S}} \label{eq:9.152}
\end{align}
entscheidet darüber, ob die Anzahl von Stützstellen in Durchlass oder Sperrbereich erhöht werden muss. Gegebenenfalls wird die Anzahl der Stützstellen angepasst und der Remez-Austauschalgorithmus erneut durchgeführt. Als Ergebnis liegt ein Amplitudengang \(A(Ω)\) vor, der aus Tschebyscheff-Polynomen besteht und das Toleranzschema erfüllt.

\subsubsection{Bestimmung von Impulsantwort und Übertragungsfunktion}\label{sec:bestimmung-von-impulsantwort-und-uxfcbertragungsfunktion}

Aus der Darstellung des Amplitudengangs \(A(Ω)\) als Tschebyscheff-Polynom kann eine Realisierungsform abgeleitet werden. Der Amplitudengang \(A(Ω)\) wird dazu über komplexe Exponentialfunktionen dargestellt als
\begin{align}
A(\Omega) &= \sum\limits_{k = 0}^{N}{a_{k}  \cos^{k}(\Omega)} = a_{0} + \sum\limits_{k = 1}^{N}\left( \frac{1}{2}  \left( e^{j  \Omega} + e^{- j  \Omega} \right) \right)^{k} \label{eq:9.153}
\end{align}
Es entstehen Summenden mit komplexen Exponentialfunktionen und Argumenten von -- jΩN \ldots{} jΩN. Durch Umsortieren ergibt sich die Darstellung
\begin{align}
A(\Omega) &= \sum\limits_{k = - N}^{N}{g_{N}[ k]  e^{- j  \Omega  k}} \label{eq:9.154}
\end{align}
Sie kann als Fourier-Transformierte der beidseitig begrenzenten Signalfolge \(g_{N}\){[}k{]} interpretiert werden. Die Signalfolge \(g_{N}\){[}k{]} ist die Impulsantwort eines Filters mit dem Amplitudengang \(A(Ω)\). Das Filter ist jedoch nicht kausal. Durch eine Verschiebung der Impulsantwort um N Takte nach rechts ergibt sich die gesuchte Impulsantwort des zu implementierenden Filters.
\begin{align}
g[ k] &= g_{N}[ k - N] \label{eq:9.155}
\end{align}
Aus der Impulsantwort errechnet sich die Übertragungsfunktion \(G(z)\) zu
\begin{align}
G(z) &= \sum\limits_{k = 0}^{\SI{2}{\newton}}{g[ k]  z^{- k}} \label{eq:9.156}
\end{align}
Damit ist das spezifizierte Filter in Zeit-, und Frequenz- und z-Bereich beschrieben.

\subsubsection{Ablauf des Remez-Filterentwurf als Flussdiagramm}\label{sec:ablauf-des-remez-filterentwurf-als-flussdiagramm}

Das Flussdiagramm in Bild \ref{fig:9.42} fasst die Vorgehensweise für den Remez-Filterentwurf zusammen.

\begin{figure}
  \centering
  \includesvg{new media/10_42_flussdiagramm_zum_remezfiltere.svg}
  \caption{Flussdiagramm zum Remez-Filterentwurf}
  \label{fig:9.42}
\end{figure}

Der hier gezeigte Ablauf ist in Programmbibliotheken hinterlegt, mit denen ein numerischer Filterentwurf schnell und unkompliziert umgesetzt werden kann. Dabei können Neben Tiefpässen auch andere Filtertypen umgesetzt werden, wie das folgende Beispiel zeigt.

\subsubsection{\texorpdfstring{Anwendungsbeispiel Equiripple-Hochpass-Filter }{Anwendungsbeispiel Equiripple-Hochpass-Filter }}\label{anwendungsbeispiel-equiripple-hochpass-filter}

Als Anwendungsbeispiel für den Remez-Entwurf wird ein Hochpass-Filter entworfen, das dem Toleranzschema in Bild \ref{fig:9.43} genügen soll.

\begin{figure}
  \centering
  \includegraphics[alt={Hochpass-Toleranzschema für das Anwendungsbeispiel}]{new media/10_43_hochpasstoleranzschema_fr_das.png}
  \caption{Hochpass-Toleranzschema für das Anwendungsbeispiel}
  \label{fig:9.43}
\end{figure}

Der Sperrbereich liegt in Bereich 0 ≤ Ω ≤ \(Ω_{S} = 0.2\)·π. In dem Bereich weist die Zielfunktion den Wert 0 auf. Die Zielfunktion wird von einem Toleranzbereich von ± \(δ_{S}\) = ± 0.1 umschlossen. Der Durchlassbereich umfasst den Bereich 0.3·\(π = Ω_{D}\) ≤ Ω ≤ π und weist ein Toleranzband von ± \(δ_{D}\) = ± 0.29 auf. Die beiden Fehlergrenzen \(δ_{D}\) und \(δ_{S}\) führen zu einer Gewichtung \(W(Ω)\) des Fehlers.
\begin{align}
W(\Omega) = \left\{ \begin{aligned}
 & \frac{\delta_{D}}{\delta_{S}} = 2.9für0 \leq \Omega \leq \Omega_{S} \\
 & 1für\Omega_{D} \leq \Omega \leq \pi
\end{aligned} \right.\ \label{eq:9.157}
\end{align}
Die erforderlichen Ordnung N wird geschätzt zu
\begin{align}
N \approx \frac{- 10  \log_{10}\left( \delta_{D}  \delta_{S} \right) - 13}{2  2.324  \left( \Omega_{D} - \Omega_{S} \right)} \label{eq:9.158}
\end{align}
Es ergibt sich eine Filterordnung von \(N = 9\). Der Remez-Entwurf führt zu einem Filter, dessen Frequenzgang in Bild \ref{fig:9.44} dargestellt ist.

\begin{figure}
  \centering
  \includegraphics[alt={Frequenzgang eines Equiripple-Hochpass-Filters der Ordnung \(N = 9\)}]{new media/10_44_frequenzgang_eines_equirippleh.png}
  \caption{Frequenzgang eines Equiripple-Hochpass-Filters der Ordnung \(N = 9\)}
  \label{fig:9.44}
\end{figure}

Der berechnete Frequenzgang stimmt aber nicht mit dem spezifizierten Toleranzschema überein. Der verwendete Algorithmus führt zu linearphasigen, symmetrische FIR-Filtern, je nach Länge entstehen verschiedene Filter der Typen I oder II:

\begin{itemize}
\item
  Typ I: Filterordnung gerade
\item
  Typ II: Filterordnung ungerade
\end{itemize}

Beide Filter haben eine symmetrische Impulsantwort, Filter von Typ II haben aber eine feste Nullstelle bei der halben Abtastfrequenz \(ω_{A}\)/2. Für einen Hochpass sind sie damit nicht geeignet. Für Hochpässe werden deshalb immer gerade Filterordnungen verwendet. Eine Erhöhung der Filterordnung auf \(N = 10\) führt zu dem in Bild \ref{fig:9.45} gezeigten Frequenzgang.

\begin{figure}
  \centering
  \includegraphics[alt={Frequenzgang eines Equiripple-Hochpass-Filters der Ordnung \(N = 10\)}]{new media/10_45_frequenzgang_eines_equirippleh.png}
  \caption{Frequenzgang eines Equiripple-Hochpass-Filters der Ordnung \(N = 10\)}
  \label{fig:9.45}
\end{figure}

Durch das Erhöhen der Filterordnung auf einen geraden Wert erfüllt der Frequenzgang das Toleranzschema. Aufgrund der Symmetrie der Filterkoeffizienten ergibt sich ein Filter mit linearer Phase.

\section{Vergleich rekursiver und nicht rekursiver Filter}\label{sec:9.5}

In den vorangegangenen Abschnitten sind einige Entwurfsmethoden für zeitdiskrete Filter beschrieben. Die unterschiedlichen Filter und Entwurfsmethoden eignen sich für unterschiedliche Aufgabenstellungen. Im Folgenden werden die Filter gegenübergestellt.

\subsubsection{Rekursive Filter (IIR-Filter)}\label{sec:rekursive-filter-iir-filter}

IIR-Filter basieren auf Filterentwürfen für zeitkontinuierliche Filter, diese Entwürfe sind standardisiert. Ausgehend von der Vorgabe eines Toleranzschemas erfolgt der Entwurf des zeitkontinuierlichen Filters nach einer klaren Prozedur. Anschließend kann das zeitkontinuierliche Filter mit Hilfe der impulsinvarianten Methode oder der bilinearen Transformation in ein zeitdiskretes Filter überführt werden. Der Filterentwurf ist damit geschlossen vorgegeben, es sind keine iterativen Verfahren zur Bestimmung der Filterkoeffizienten notwendig. Andererseits können aber auch keine spezifischen Modifikationen durchgeführt werden.

Das Verhältnis von Steilheit zur Filterordnung ist bei IIR-Filtern grundsätzlich höher als bei den beschriebenen FIR-Filtern.

Rekursive Filter ergeben sich aus der Spezifikation eines Amplitudengangs, der Phasengang bleibt während des gesamten Entwurfsverfahrens weitgehend unberücksichtigt. Rekursive Filter eigenen sich deshalb nicht dazu, linearphasige Filter oder Filter mit fest vorgegebenem Phasengang zu entwickeln.

\subsubsection{Nicht rekursive Filter (FIR-Filter)}\label{sec:nicht-rekursive-filter-fir-filter}

FIR-Filter werden nicht auf Basis eines Toleranzschemas entwickelt, sie besitzen keine geschlossenen Entwurfsgleichungen. Dadurch kann eine Iteration des Entwurfes notwendig sein, um die Vorgaben eines Toleranzschemas zu erfüllen. Bei dem Einsatz von Software für den Filterentwurf ist diese Iteration jedoch unkompliziert möglich. Durch den Einsatz von Optimierungskriterien im Entwurfsprozess wie zum Beispiel beim Remez-Entwurf können sehr spezifische Filter entwickelt werden.

Das Verhältnis von Steilheit zu Filterordnung ist grundsätzlich schlechter als bei IIR-Filtern, durch den Einsatz von Dezimierungsstufen und Cascaded-Interpolated-Comb-Strukturen (CIC-Filter) kann die Effizienz von FIR-Filtern jedoch gesteigert werden.

Wesentlicher Vorteil von FIR-Filtern ist ihr linearer Phasengang.

\subsubsection{Tabellarischer Vergleich}\label{sec:tabellarischer-vergleich}

Tabelle \ref{tab:9.9} fasst den Vergleich von IIR- und FIR-Filtern tabellarisch zusammen.

\begin{table}
  \centering
  \caption{Tabellarischer Vergleich von IIR- und FIR-Filtern}
  \label{tab:9.9}
  \begin{tabular}{lll}
    \toprule
    Eigenschaft & IIR-Filter & FIR-Filter \\
    \midrule
    Entwurfsverfahren & 
Geschlossenes 
Entwurfsverfahren
 & 
Approximations- 
verfahren \\
    Einfluss auf den Filterentwurf & 
Wenig Freiheitsgrade, 
geschlossenes Entwurfsverfahren
 & Optimierungskriterien definierbar, zum Beispiel Remez-Entwurf \\
    Effizienz & 
Hohe Effizienz bei Filtern, 
weitgehend ohne Phasenvorgaben
 & 
Grundsätzlich schlechtere 
Effizienz als bei IIR-Filtern \\
    Lineare Phase & Nur in kleinen Frequenzintervallen annähernd lineare Phase & 
Filter mit linearer Phase 
realisierbar \\
    \bottomrule
  \end{tabular}
\end{table}
\begin{application}{Entwurf Digitaler Filter}
  Im Online-Portal \emph{Systemtheorie Online} verdeutlicht die Applikation \emph{Entwurf Digitaler Filter} den Filterentwurf für IIR- und FIR-Filter.
  \tcblower
  \url{h-ka.de/eit/digitale-signalverarbeitung/übungen/kapitel2}
\end{application}
\section{Strukturen und Blockschaltbilder digitaler Systeme}\label{sec:9.6}

Die in diesem Kapitel entworfenen Systeme lassen sich über Differenzengleichungen und Übertragungsfunktionen charakterisieren. Die Realisierung zeitdiskreter Systeme kann bei gleicher Übertragungsfunktion aber auf unterschiedliche Art erfolgen. Zur Beschreibung der Realisierungsvarianten werden Blockschaltbilder verwendet. Sie sind Gegenstand dieses Abschnittes.

Nach einer Einführung in Blockschaltbilder zeitdiskreter Systeme werden drei grundlegende Strukturen zeitdiskreter Systeme dargestellt und diskutiert:

\begin{itemize}
\item
  Direktstruktur
\item
  Kaskadenstruktur
\item
  Parallelstruktur
\end{itemize}

Die Strukturen stellen symbolisch den Aufbau von Schaltungsstrukturen oder Algorithmen dar. Sie lassen sich ineinander überführen und sind theoretisch gleichwertig. Durch die unterschiedliche Realisierung ergeben sich praktisch jedoch Unterschiede. Die Anzahl der benötigten Speicherzellen und der erforderlichen Rechenoperationen variiert. Damit ist auch der Realisierungsaufwand der Strukturen unterschiedlich. Außerdem führt bei der sogenannten Fixed-Point-Darstellung der Signale und Koeffizienten die Quantisierung zu Abweichungen. Dadurch wird das System nicht ideal abgebildet. Es entstehen Quantisierungsfehler, die im Extremfall sogar zur Instabilität von Systemen führen können. Die Auswirkung der Quantisierungseffekte ist abhängig von der Systemrealisierung.

\subsection{Einführung in Blockschaltbilder zeitdiskreter Systeme}\label{sec:einfuxfchrung-in-blockschaltbilder-zeitdiskreter-systeme}

Zeitdiskrete Systeme können wie zeitkontinuierliche Systeme als Blockschaltbild dargestellt werden. Als Einführungsbeispiel dient das Blockschaltbild der Differenzengleichung
\begin{align}
y[ k] &= u[ k] + a  y[ k - 1] \label{eq:9.159}
\end{align}
Wie bei zeitkontinuierlichen Systemen wird eine Verstärkung mit einem Faktor an der Linie dargestellt, Summationspunkte über Kreise und einzelne Übertragungsglieder in einem Rechteck. Pfeile geben die Flussrichtung des Signals an. Das Blockschaltbild ist in Bild \ref{fig:9.46} dargestellt.

\begin{figure}
  \centering
  \includegraphics[alt={Darstellung eines zeitdiskreten Systems als Blockschaltbild}]{new media/10_46_darstellung_eines_zeitdiskrete.png}
  \caption{Darstellung eines zeitdiskreten Systems als Blockschaltbild}
  \label{fig:9.46}
\end{figure}

Zu dem Eingangssignal \(u[k]\) wird ein um einen Takt verzögertes und mit dem Faktor a multipliziertes Ausgangssignal addiert. Das Ergebnis ist das Ausgangssignal \(y[k]\). Das System beschreibt demnach obige Differenzengleichung.

\subsection{Direktstrukturen}\label{sec:direktstrukturen}

Die unterschiedlichen Strukturen von Blockschaltbildern ergeben sich aus unterschiedlichen Darstellungen von Übertragungsfunktionen. Generell werden bei Systemrealisierungen nur Verzögerungsglieder mit der z-Transformierten \(z^{-1}\) eingesetzt. Andernfalls ließe sich das System wegen der nicht gegebenen Kausalität nicht realisieren. Die Direktstruktur geht deshalb von einer Übertragungsfunktion aus, die als gebrochen rationale Funktion mit negativen Exponenten dargestellt wird.
\begin{align}
G(z) &= \frac{Y(z)}{U(z)} = \frac{\sum\limits_{m = 0}^{M}{d_{m}  z^{- m}}}{\sum\limits_{n = 0}^{N}{c_{n}  z^{- n}}} = \frac{d_{0} + d_{1}  z^{- 1} + ... + d_{M}  z^{- M}}{c_{0} + c_{1}  z^{- 1} + ... + c_{N}  z^{- N}} \label{eq:9.160}
\end{align}
Dabei kann der Koeffizient \(c_{0} = 1\) gewählt werden. Ist das nicht der Fall, wird der gesamte Bruch durch \(c_{0}\) gekürzt. Mit dieser Voraussetzung errechnet sich die Ausgangsgröße im Bildbereich zu
\begin{align}
Y(z) &= \left( d_{0} + d_{1}  z^{- 1} + ... + d_{M}  z^{- M} \right)  U(z) - \left( c_{1}  z^{- 1} + ... + c_{N}  z^{- N} \right)  Y(z) \label{eq:9.161}
\end{align}
Diese Gleichung führt zu der Direktstruktur 1 für zeitdiskrete Systeme, die in Bild \ref{fig:9.47} dargestellt ist.

\begin{figure}
  \centering
  \includegraphics[alt={Direktstruktur 1 eines zeitdiskreten Systems}]{new media/10_47_direktstruktur_1_eines_zeitdis.png}
  \caption{Direktstruktur 1 eines zeitdiskreten Systems}
  \label{fig:9.47}
\end{figure}

Ein Vergleich von Gleichung \eqref{eq:9.161} und des Blockschaltbildes zeigt, dass die Direktstruktur 1 und die Gleichung dieselbe Übertragungsfunktion besitzen. Bei der Darstellung dieser Struktur werden Zählergrad M und Nennergrad N gleichgesetzt. Ist das nicht der Fall, werden die entsprechenden Koeffizienten \(d_{m}\) beziehungsweise \(c_{n}\) zu null gesetzt und das Blockdiagramm gegebenenfalls vereinfacht.

\begin{example}{Implementierung eines Systems als Direktstruktur 1}
\label{beispiel-implementierung-eines-systems-als-direktstruktur-1}

Als Beispiel wird ein System mit der Übertragungsfunktion \(G(z)\) in Direktstruktur 1 dargestellt.
\begin{align}
G(z) &= \frac{z^{- 1} - 1.25  z^{- 2} + 0.25  z^{- 3}}{1 - 1.5  z^{- 1} + z^{- 2} - 0.25  z^{- 3}} \label{eq:9.162}
\end{align}
Einsetzen der Koeffizienten aus der gebrochen rationalen Funktion führt zu dem in Bild \ref{fig:9.48} dargestellten Blockschaltbild.

\begin{fixedfigure}
  \includegraphics[alt={Beispiel für eine Systemdarstellung in Direktstruktur 1}]{new media/10_48_beispiel_fr_eine_systemdarstel.png}
  \captionof{figure}{Beispiel für eine Systemdarstellung in Direktstruktur 1}
  \label{fig:9.48}
\end{fixedfigure}

Für die Realisierung werden 6 Speicher, 6 Multiplikationen und 5 Additionen benötigt.

\end{example}

Die Darstellung in Direktstruktur 1 erfordert im Allgemeinen N + M Speicherelemente, weil die Verzögerungsglieder für Eingangssignal und Ausgangssignal separat verwendet werden. Um Speicherelemente einzusparen, kann die Differenzengleichung umgeformt werden zu
\begin{align}
Y(z) &= d_{0}  U(z) + \left( d_{1}  X(z) - c_{1}  Y(z) \right)  z^{- 1} + ... + \left( d_{N}  U(z) - c_{N}  Y(z) \right)  z^{- N} \label{eq:9.163}
\end{align}
Diese Gleichung führt zu einer Realisierungsform mit minimaler Anzahl von Speicherelementen. Die in Gleichung \eqref{eq:9.163} beschriebene Struktur wird transponierte Direktstruktur 2 genannt und ist in Bild \ref{fig:9.49} dargestellt.

\begin{figure}
  \centering
  \includegraphics[alt={Transponierte Direktstruktur 2 eines zeitdiskreten Systems}]{new media/10_49_transponierte_direktstruktur_2.png}
  \caption{Transponierte Direktstruktur 2 eines zeitdiskreten Systems}
  \label{fig:9.49}
\end{figure}

\begin{example}{Implementierung eines Systems als transponierte Direktstruktur 2}
\label{beispiel-implementierung-eines-systems-als-transponierte-direktstruktur-2}

Als Beispiel wird ein System mit der Übertragungsfunktion \(G(z)\) in transponierter Direktstruktur 2 dargestellt.
\begin{align}
G(z) &= \frac{z^{- 1} - 1.25  z^{- 2} + 0.25  z^{- 3}}{1 - 1.5  z^{- 1} + z^{- 2} - 0.25  z^{- 3}} \label{eq:9.164}
\end{align}
Einsetzen der Koeffizienten führt zu dem in Bild \ref{fig:9.50} dargestellten Signalflussgraphen.

\begin{fixedfigure}
  \includegraphics[alt={Beispiel für eine Systemdarstellung in transponierter Direktstruktur 2}]{new media/10_50_beispiel_fr_eine_systemdarstel.png}
  \captionof{figure}{Beispiel für eine Systemdarstellung in transponierter Direktstruktur 2}
  \label{fig:9.50}
\end{fixedfigure}

Für die Realisierung werden nur noch 3 Speicher benötigt, die Anzahl der Multiplikationen und Additionen bleibt gleich.

\end{example}

\subsection{Kaskadenstrukturen}\label{sec:kaskadenstrukturen}

Bei der Darstellung des Systems in Direktstruktur 1 oder 2 wird die gebrochen rationale Funktion als Ausgangspunkt betrachtet. Ändert sich ein Koeffizient der Übertragungsfunktion zum Beispiel wegen Quantisierungsfehlern, ändert sich die gesamte Übertragungsfunktion, weil sämtliche Koeffizienten in die Berechnung der Pole eingehen. Alternativ zu den Direktstrukturen werden deshalb Kaskadenstrukturen aufgebaut, die aus dem Produkt einzelner Teilübertragungsfunktionen zweiter Ordnung gebildet werden. Die Übertragungsfunktion \(G(z)\) wird dabei in das Produkt von Teil-Übertragungsfunktionen \(G_{n}(z)\) zerlegt.
\begin{align}
G(z) &= G_{1}(z)  G_{2}(z)  ... = \frac{d_{01} + d_{11}  z^{- 1} + d_{21}  z^{- 2}}{c_{01} + c_{11}  z^{- 1} + c_{21}  z^{- 2}}  \frac{d_{02} + d_{12}  z^{- 1} + d_{22}  z^{- 2}}{c_{02} + c_{12}  z^{- 1} + c_{22}  z^{- 2}}  ... \label{eq:9.165}
\end{align}
Um Systeme mit reellen Koeffizienten zu erhalten, werden Systeme zweiter Ordnung verwendet. Jeweils zwei reelle oder konjugiert komplexe Pole und Nullstellen führen zu einem Nenner- oder Zählerpolynom zweiter Ordnung mit reellen Koeffizienten. Bei Systemen ungerader Ordnung wird die Struktur mit einem Teilsystem erster Ordnung ergänzt. Die einzelnen Teilsysteme werden als Biquad-Strukturen bezeichnet und in der transponierten Direktstruktur 2 dargestellt.

\begin{figure}
  \centering
  \includegraphics[alt={Biquad-Element}]{new media/10_51_biquadelement.png}
  \caption{Biquad-Element}
\end{figure}

Der Vorteil der Biquads liegt darin, dass Änderungen der Koeffizienten sich nur auf das Pol- oder Nullstellenpaar auswirken, das mit ihm dargestellt wird. Dadurch wird das Gesamtsystem numerisch stabiler und übersichtlicher.

\begin{example}{Implementierung eines Systems als Kaskadenstruktur}
\label{beispiel-implementierung-eines-systems-als-kaskadenstruktur}

Als Beispiel wird wieder ein System mit der Übertragungsfunktion \(G(z)\) dargestellt. Die Übertragungsfunktion kann in zwei Faktoren zerlegt werden.
\begin{align}
G(z) &= \frac{z^{- 1} - z^{- 2}}{1 - 1.2  z^{- 1} + 0.7  z^{- 2} - 0.1  z^{- 3}} = \frac{1 + 1.25z^{- 1} + 0.25  z^{- 2}}{1 - z^{- 1} + 0.5  z^{- 2}}  \frac{1}{1 - 0.5  z^{- 1}} \label{eq:9.166}
\end{align}
Die Darstellung führt zu dem in Blockschaltbild mit Biquads, das in Bild \ref{fig:9.52} dargestellt ist.

\begin{fixedfigure}
  \includegraphics[alt={Beispiel zur Kaskadenstruktur mit Biquads}]{new media/10_52_beispiel_zur_kaskadenstruktur.png}
  \captionof{figure}{Beispiel zur Kaskadenstruktur mit Biquads}
  \label{fig:9.52}
\end{fixedfigure}

Wie bei der transponierten Direktstruktur 2 werden 3 Speicher benötigt, auch die Anzahl der Multiplikationen und Additionen bleibt gleich.

\end{example}

\subsection{Parallelstrukturen}\label{sec:parallelstrukturen}

Alternativ zur Kaskadenstruktur kann ein Filter als Parallelstruktur realisiert werden. Dazu wird seine Übertragungsfunktion in Partialbrüche zerlegt und Partialbrüche mit konjugiert komplexen Polpaaren zu einem Biquad zusammengefasst. Es entsteht eine Übertragungsfunktion der Form
\begin{align}
G(z) &= G_{1}(z) + G_{2}(z)  ... = \frac{d_{01} + d_{11}  z^{- 1} + d_{21}  z^{- 2}}{c_{01} + c_{11}  z^{- 1} + c_{21}  z^{- 2}} + \frac{d_{02} + d_{12}  z^{- 1} + d_{22}  z^{- 2}}{c_{02} + c_{12}  z^{- 1} + c_{22}  z^{- 2}}  ... \label{eq:9.167}
\end{align}
Für ein System vierter Ordnung ergibt sich damit eine Parallelschaltung zweier Biquads, wie sie in Bild \ref{fig:9.53} dargestellt ist.

\begin{figure}
  \centering
  \includegraphics[alt={Parallelstruktur mit Biquads}]{new media/10_53_parallelstruktur_mit_biquads.png}
  \caption{Parallelstruktur mit Biquads}
  \label{fig:9.53}
\end{figure}

Alle Biquads erhalten das Eingangssignal zum gleichen Zeitpunkt. Dadurch wird das parallele Rechnen der Systemantwort möglich.

\begin{example}{Implementierung eines Systems als Parallelstruktur}
\label{beispiel-implementierung-eines-systems-als-parallelstruktur}

Als Beispiel wird wieder ein System mit der Übertragungsfunktion \(G(z)\) dargestellt. Die Übertragungsfunktion kann in zwei Summanden mit reellen Koeffizienten zerlegt werden.
\begin{align}
G(z) &= \frac{z^{- 1} - z^{- 2}}{1 - 1.2  z^{- 1} + 0.7  z^{- 2} - 0.1  z^{- 3}} = \frac{1.5  z^{- 1} + 1  z^{- 2}}{1 - z^{- 1} + 0.5  z^{- 2}} + \frac{\text{-0.5}}{1 - 0.5  z^{- 1}} \label{eq:9.168}
\end{align}
Die Darstellung führt zu dem Blockschaltbild mit Biquads, das in Bild \ref{fig:9.54} dargestellt ist.

\begin{fixedfigure}
  \includegraphics[alt={Beispiel zur Parallelstruktur mit Biquads}]{new media/10_54_beispiel_zur_parallelstruktur.png}
  \captionof{figure}{Beispiel zur Parallelstruktur mit Biquads}
  \label{fig:9.54}
\end{fixedfigure}

Auch bei dieser Realisierungsform werden 3 Speicher benötigt und die Anzahl der Multiplikationen und Additionen ändert sich gegenüber der Direktstruktur 2 nicht.

\end{example}

\subsection{Transversale Strukturen}\label{sec:transversale-strukturen}

Die in den vorangegangenen Abschnitten dargestellten Strukturen eignen sich für Systeme mit Rückkopplung und damit für Systeme mit unendlicher Impulsantwort (IIR-Systeme). Bei FIR-Systemen existieren keine Rückkopplungen und die Impulsantwort ist endlich. Aus der Übertragungsfunktion ist direkt ersichtlich, dass bei FIR-Systemen nur das Eingangssignal verzögert werden muss.
\begin{align}
G(z) &= \frac{Y(z)}{U(z)} = \sum\limits_{m = 0}^{M}{d_{m}  z^{- m}} \label{eq:9.169}
\end{align}
Es ergibt sich eine vergleichsweise einfache Realisierung, die als transversale Struktur bezeichnet wird. Sie ist in Bild \ref{fig:9.55} dargestellt.

\begin{figure}
  \centering
  \includegraphics[alt={Transversale Struktur}]{new media/10_55_transversale_struktur.png}
  \caption{Transversale Struktur}
  \label{fig:9.55}
\end{figure}

Das Eingangssignal durchläuft die gesamte Signalkette, sodass Fehler durch gerundete Korrekturen bei jedem Summanden nur einmal vorkommen können und keine Fehlerkumulation auftritt.

\begin{example}{Implementierung eines Systems als transversale Struktur}
\label{beispiel-implementierung-eines-systems-als-transversale-struktur}

Als Beispiel wird der gleitende Mittelwert aus 4 Gliedern dargestellt. Er hat die Übertragungsfunktion
\begin{align}
G(z) &= \frac{z^{3} + z^{2} + z + 1}{4  z^{3}} = \frac{1}{4}  \left( 1 + z^{- 1} + z^{- 2} + z^{- 3} \right) \label{eq:9.170}
\end{align}
Es ergibt sich die in Bild \ref{fig:9.56} dargestellte Transversalstruktur.

\begin{fixedfigure}
  \includegraphics[alt={Beispiel für eine transversale Struktur}]{new media/10_56_beispiel_fr_eine_transversale.png}
  \captionof{figure}{Beispiel für eine transversale Struktur}
  \label{fig:9.56}
\end{fixedfigure}

\end{example}

\subsection{\texorpdfstring{Fixed-Point-Darstellung }{Fixed-Point-Darstellung }}\label{fixed-point-darstellung}

\textless\textless\textless{} wird später ergänzt \textgreater\textgreater\textgreater{}

\subsection{Bewertung der Realisierungsstrukturen unter Berücksichtigung der Fixed-Point-Darstellung}\label{sec:bewertung-der-realisierungsstrukturen-unter-beruxfccksichtigung-der-fixed-point-darstellung}
\begin{application}{Quantisierung von Filterkoeffizienten}
  Im Online-Portal \emph{Systemtheorie Online} verdeutlicht die Applikation \emph{Quantisierung von Filterkoeffizienten} die Auswirkung einer Fixed-Point-Darstellung bei Filtern.
  \tcblower
  \url{h-ka.de/eit/digitale-signalverarbeitung/übungen/kapitel2}
\end{application}
